\documentclass[a4paper,10pt]{article}
\usepackage{eumat}

\begin{document}
\begin{eulernotebook}
\eulerheading{Menggambar Grafik 2D dengan EMT}
\begin{eulercomment}
Notebook ini menjelaskan tentang cara menggambar berbagaikurva dan
grafik 2D dengan software EMT. EMT menyediakan fungsi plot2d() untuk
menggambar berbagai kurva dan grafik dua dimensi (2D).

\begin{eulercomment}
\eulerheading{Soal latihan ada di bawah}
\begin{eulercomment}
\end{eulercomment}
\eulersubheading{Dasar-Dasar Plot}
\begin{eulercomment}
Ada fungsi-fungsi dasar untuk menggambar plot. Terdapat koordinat
layar, yang selalu berkisar dari 0 sampai 1024 pada setiap sumbu,
tidak peduli apakah layarnya persegi atau tidak. Selain itu ada
koordinat plot, yang dapat diatur dengan setplot(). Pemetaan antara
kedua koordinat tersebut bergantung pada jendela plot yang sedang
aktif. Misalnya, shrinkwindow() secara default menyisakan ruang untuk
label sumbu dan judul plot.

Pada contoh ini, kita hanya menggambar beberapa garis acak dengan
berbagai warna. Untuk rincian fungsi-fungsi tersebut, pelajari
fungsi-fungsi inti EMT.
\end{eulercomment}
\begin{eulerprompt}
>clg; // hapus layar
>window(0,0,1024,1024); // gunakan seluruh jendela
>setplot(0,1,0,1); // atur koordinat plot
>hold on; // mulai mode timpa
>n=100; X=random(n,2); Y=random(n,2);  // bangkitkan titik acak
>colors=rgb(random(n),random(n),random(n)); // bangkitkan warna acak
>loop 1 to n; color(colors[#]); plot(X[#],Y[#]); end; // gambar
>hold off; // akhiri mode timpa
>insimg; // sisipkan ke notebook
\end{eulerprompt}
\eulerimg{13}{images/EMT4Plot2D_Indonesia-001.png}
\begin{eulerprompt}
>reset;
\end{eulerprompt}
\begin{eulercomment}
Grafik perlu ditahan (hold), karena perintah plot() akan menghapus jendela plot.

Untuk menghapus semua yang telah kita lakukan, kita gunakan reset().

Untuk menampilkan gambar hasil plot di layar notebook, perintah plot2d() dapat diakhiri
dengan titik dua (:). Cara lain adalah perintah plot2d() diakhiri dengan titik koma (;),
kemudian menggunakan perintah insimg() untuk menampilkan gambar hasil plot.

Sebagai contoh lain, kita menggambar sebuah plot sebagai sisipan (inset) di dalam plot
lain. Ini dilakukan dengan mendefinisikan jendela plot yang lebih kecil. Perhatikan
bahwa jendela ini tidak menyediakan ruang untuk label sumbu di luar jendela plot,
sehingga kita perlu menambahkan margin seperlunya. Perhatikan juga bahwa kita
menyimpan dan memulihkan jendela penuh, serta menahan (hold) plot yang sedang aktif
selagi menggambar sisipannya.
\end{eulercomment}
\begin{eulerprompt}
>plot2d("x^3-x");
>xw=200; yw=100; ww=300; hw=300;
>ow=window();
>window(xw,yw,xw+ww,yw+hw);
>hold on;
>barclear(xw-50,yw-10,ww+60,ww+60);
>plot2d("x^4-x",grid=6):
\end{eulerprompt}
\eulerimg{13}{images/EMT4Plot2D_Indonesia-002.png}
\begin{eulerprompt}
>hold off;
>window(ow);
\end{eulerprompt}
\begin{eulercomment}
Plot dengan beberapa gambar (multi-figure) dapat dibuat dengan cara yang sama. Untuk
ini tersedia fungsi utilitas figure().

\end{eulercomment}
\eulersubheading{Aspek Plot}
\begin{eulercomment}
Plot default menggunakan jendela plot persegi. Anda dapat mengubahnya dengan fungsi
aspect(). Jangan lupa mengatur ulang aspek tersebut setelahnya. Anda juga dapat
mengubah nilai default ini melalui menu "Set Aspect" ke rasio aspek tertentu atau ke
ukuran jendela grafik saat ini.

Anda juga dapat mengubahnya hanya untuk satu plot. Untuk itu, ukuran area plot saat
ini diubah, dan jendela diatur sehingga label memiliki cukup ruang.
\end{eulercomment}
\begin{eulerprompt}
>aspect(2); // rasio panjang dan lebar 2:1
>plot2d(["sin(x)","cos(x)"],0,2pi):
\end{eulerprompt}
\eulerimg{13}{images/EMT4Plot2D_Indonesia-003.png}
\begin{eulerprompt}
>aspect();
>reset;
\end{eulerprompt}
\begin{eulercomment}
Fungsi reset() mengembalikan pengaturan default plot, termasuk rasio aspek.\\
\begin{eulercomment}
\eulerheading{Plot 2D di Euler}
\begin{eulercomment}
EMT Math Toolbox memiliki plot 2D, baik untuk data maupun fungsi. EMT menggunakan
fungsi plot2d. Fungsi ini dapat menggambar fungsi dan data.

Dimungkinkan juga menggambar plot di Maxima menggunakan Gnuplot atau di Python
menggunakan Math Plot Lib.

Euler dapat membuat plot 2D dari

- ekspresi\\
- fungsi, variabel, atau kurva parametrik,\\
- vektor nilai-nilai x-y,\\
- awan titik pada bidang,\\
- kurva implisit dengan level atau daerah level.\\
- fungsi kompleks

Gaya plot meliputi berbagai gaya untuk garis dan titik, plot batang, dan plot
berarsir.\\
\begin{eulercomment}
\eulerheading{Plot Ekspresi atau Variabel}
\begin{eulercomment}
Satu ekspresi dalam "x" (misalnya "4*x\textasciicircum{}2") atau nama sebuah fungsi (misalnya "f")
menghasilkan grafik fungsi tersebut.

Berikut contoh paling dasar, yang menggunakan rentang default dan mengatur rentang y
yang sesuai agar cocok dengan plot fungsi.

Catatan: Jika Anda mengakhiri baris perintah dengan titik dua ":", plot akan
disisipkan ke jendela teks. Jika tidak, tekan TAB untuk melihat plot jika jendela plot
tertutup.
\end{eulercomment}
\begin{eulerprompt}
>plot2d("x^2"):
\end{eulerprompt}
\eulerimg{13}{images/EMT4Plot2D_Indonesia-004.png}
\begin{eulerprompt}
>aspect(1.5); plot2d("x^3-x"):
\end{eulerprompt}
\eulerimg{17}{images/EMT4Plot2D_Indonesia-005.png}
\begin{eulerprompt}
>a:=5.6; plot2d("exp(-a*x^2)/a"); insimg(30); // menampilkan gambar hasil plot setinggi 25 baris
\end{eulerprompt}
\eulerimg{17}{images/EMT4Plot2D_Indonesia-006.png}
\begin{eulercomment}
Dari beberapa contoh sebelumnya Anda dapat melihat bahwa aslinya gambar plot
menggunakan sumbu X dengan rentang nilai dari -2 sampai dengan 2. Untuk mengubah
rentang nilai X dan Y, Anda dapat menambahkan nilai-nilai batas X (dan Y) di
belakang ekspresi yang digambar.

Rentang plot diatur dengan parameter-parameter berikut

- a,b: rentang x (default -2,2)\\
- c,d: rentang y (default: menyesuaikan dengan nilai)\\
- r: alternatif berupa jari-jari di sekitar pusat plot\\
- cx,cy: koordinat pusat plot (default 0,0)
\end{eulercomment}
\begin{eulerprompt}
>plot2d("x^3-x",-1,2):
\end{eulerprompt}
\eulerimg{17}{images/EMT4Plot2D_Indonesia-007.png}
\begin{eulerprompt}
>plot2d("sin(x)",-2*pi,2*pi): // plot sin(x) pada interval [-2pi, 2pi]
\end{eulerprompt}
\eulerimg{17}{images/EMT4Plot2D_Indonesia-008.png}
\begin{eulerprompt}
>plot2d("cos(x)","sin(3*x)",xmin=0,xmax=2pi):
\end{eulerprompt}
\eulerimg{17}{images/EMT4Plot2D_Indonesia-009.png}
\begin{eulercomment}
Alternatif selain titik dua adalah perintah insimg(lines), yang menyisipkan plot
dengan menempati sejumlah baris teks tertentu.

Pada opsi, plot dapat diatur agar muncul

- di jendela terpisah yang dapat diubah ukurannya,\\
- di jendela notebook.

Gaya yang lebih beragam dapat dicapai dengan perintah plot tertentu.

Bagaimanapun, tekan tombol tabulator untuk melihat plot jika plot tersembunyi.

Untuk membagi jendela menjadi beberapa plot, gunakan perintah figure(). Pada contoh,
kita menggambar x\textasciicircum{}1 sampai x\textasciicircum{}4 ke dalam 4 bagian jendela. figure(0) mengatur ulang
jendela default.
\end{eulercomment}
\begin{eulerprompt}
>reset;
>figure(2,2); ...
>for n=1 to 4; figure(n); plot2d("x^"+n); end; ...
>figure(0):
\end{eulerprompt}
\eulerimg{13}{images/EMT4Plot2D_Indonesia-010.png}
\begin{eulercomment}
Pada plot2d(), tersedia gaya alternatif dengan grid=x. Sebagai gambaran, kita
tampilkan berbagai gaya grid dalam satu gambar (lihat di bawah untuk perintah
figure()). Gaya grid=0 tidak disertakan. Gaya ini tidak menampilkan grid maupun
bingkai.
\end{eulercomment}
\begin{eulerprompt}
>figure(3,3); ...
>for k=1:9; figure(k); plot2d("x^3-x",-2,1,grid=k); end; ...
>figure(0):
\end{eulerprompt}
\eulerimg{13}{images/EMT4Plot2D_Indonesia-011.png}
\begin{eulercomment}
Jika argumen plot2d() berupa sebuah ekspresi diikuti empat bilangan, bilangan-bilangan
tersebut adalah rentang x dan y untuk plot.

Alternatifnya, a, b, c, d dapat ditentukan sebagai parameter yang ditetapkan, yaitu
a=... dan seterusnya.

Pada contoh berikut, kita mengubah gaya grid, menambahkan label, dan menggunakan label
vertikal untuk sumbu y.
\end{eulercomment}
\begin{eulerprompt}
>aspect(1.5); plot2d("sin(x)",0,2pi,-1.2,1.2,grid=3,xl="x",yl="sin(x)"):
\end{eulerprompt}
\eulerimg{17}{images/EMT4Plot2D_Indonesia-012.png}
\begin{eulerprompt}
>plot2d("sin(x)+cos(2*x)",0,4pi):
\end{eulerprompt}
\eulerimg{17}{images/EMT4Plot2D_Indonesia-013.png}
\begin{eulercomment}
Gambar yang dihasilkan dengan menyisipkan plot ke jendela teks disimpan di direktori
yang sama dengan notebook, secara default pada subdirektori bernama "images".
Gambar-gambar ini juga digunakan oleh ekspor HTML.

Anda cukup menandai gambar mana pun dan menyalinnya ke clipboard dengan Ctrl-C. Tentu
saja, Anda juga dapat mengekspor grafik saat ini dengan fungsi-fungsi pada menu File.

Fungsi atau ekspresi pada plot2d dievaluasi secara adaptif. Untuk kecepatan lebih
tinggi, matikan plot adaptif dengan \textless{}adaptive dan tentukan banyaknya subinterval
dengan n=... Hal ini hanya diperlukan pada kasus yang jarang.
\end{eulercomment}
\begin{eulerprompt}
>plot2d("sign(x)*exp(-x^2)",-1,1,<adaptive,n=10000):
\end{eulerprompt}
\eulerimg{17}{images/EMT4Plot2D_Indonesia-014.png}
\begin{eulerprompt}
>plot2d("x^x",r=1.2,cx=1,cy=1):
\end{eulerprompt}
\eulerimg{17}{images/EMT4Plot2D_Indonesia-015.png}
\begin{eulercomment}
Perhatikan bahwa x\textasciicircum{}x tidak terdefinisi untuk x\textless{}=0. Fungsi plot2d menangkap galat ini,
dan mulai menggambar segera setelah fungsi terdefinisi. Ini berlaku untuk semua fungsi
yang mengembalikan NAN di luar daerah definisinya.
\end{eulercomment}
\begin{eulerprompt}
>plot2d("log(x)",-0.1,2):
\end{eulerprompt}
\eulerimg{17}{images/EMT4Plot2D_Indonesia-016.png}
\begin{eulercomment}
Parameter square=true (atau \textgreater{}square) memilih rentang y secara otomatis sehingga
hasilnya berupa jendela plot persegi. Perhatikan bahwa secara default, Euler
menggunakan ruang persegi di dalam jendela plot.
\end{eulercomment}
\begin{eulerprompt}
>plot2d("x^3-x",>square):
\end{eulerprompt}
\eulerimg{17}{images/EMT4Plot2D_Indonesia-017.png}
\begin{eulerprompt}
>plot2d(''integrate("sin(x)*exp(-x^2)",0,x)'',0,2): // gambar integral
\end{eulerprompt}
\eulerimg{17}{images/EMT4Plot2D_Indonesia-018.png}
\begin{eulercomment}
Jika Anda memerlukan lebih banyak ruang untuk label y, panggil shrinkwindow() dengan
parameter smaller, atau atur nilai positif untuk "smaller" pada plot2d().
\end{eulercomment}
\begin{eulerprompt}
>plot2d("gamma(x)",1,10,yl="y-values",smaller=6,<vertical):
\end{eulerprompt}
\eulerimg{17}{images/EMT4Plot2D_Indonesia-019.png}
\begin{eulercomment}
Ekspresi simbolik juga dapat digunakan, karena disimpan sebagai ekspresi string
sederhana.
\end{eulercomment}
\begin{eulerprompt}
>x=linspace(0,2pi,1000); plot2d(sin(5x),cos(7x)):
\end{eulerprompt}
\eulerimg{17}{images/EMT4Plot2D_Indonesia-020.png}
\begin{eulerprompt}
>a:=5.6; expr &= exp(-a*x^2)/a; // definisikan ekspresi
>plot2d(expr,-2,2): // gambar dari -2 sampai 2
\end{eulerprompt}
\eulerimg{17}{images/EMT4Plot2D_Indonesia-021.png}
\begin{eulerprompt}
>plot2d(expr,r=1,thickness=2): // gambar dalam persegi di sekitar (0,0)
\end{eulerprompt}
\eulerimg{17}{images/EMT4Plot2D_Indonesia-022.png}
\begin{eulerprompt}
>plot2d(&diff(expr,x),>add,style="--",color=red): // tambahkan plot lain
\end{eulerprompt}
\eulerimg{17}{images/EMT4Plot2D_Indonesia-023.png}
\begin{eulerprompt}
>plot2d(&diff(expr,x,2),a=-2,b=2,c=-2,d=1): // gambar dalam persegi panjang
\end{eulerprompt}
\eulerimg{17}{images/EMT4Plot2D_Indonesia-024.png}
\begin{eulerprompt}
>plot2d(&diff(expr,x),a=-2,b=2,>square): // pertahankan plot persegi
\end{eulerprompt}
\eulerimg{17}{images/EMT4Plot2D_Indonesia-025.png}
\begin{eulerprompt}
>plot2d("x^2",0,1,steps=1,color=red,n=10):
\end{eulerprompt}
\eulerimg{17}{images/EMT4Plot2D_Indonesia-026.png}
\begin{eulerprompt}
>plot2d("x^2",>add,steps=2,color=blue,n=10):
\end{eulerprompt}
\eulerimg{17}{images/EMT4Plot2D_Indonesia-027.png}
\eulerheading{Fungsi dengan Satu Parameter}
\begin{eulercomment}
Fungsi plot yang paling penting untuk plot bidang adalah plot2d(). Fungsi ini
diimplementasikan dalam bahasa Euler pada file "plot.e", yang dimuat saat program
dimulai.

Berikut beberapa contoh yang menggunakan fungsi. Seperti biasa di EMT, untuk fungsi
yang bekerja pada fungsi atau ekspresi lain, Anda dapat meneruskan parameter tambahan
(selain x) yang bukan variabel global ke fungsi dengan parameter titik koma atau
dengan call collection.
\end{eulercomment}
\begin{eulerprompt}
>function f(x,a) := x^2/a+a*x^2-x; // definisikan sebuah fungsi
>a=0.3; plot2d("f",0,1;a): // gambar dengan a=0.3
\end{eulerprompt}
\eulerimg{17}{images/EMT4Plot2D_Indonesia-028.png}
\begin{eulerprompt}
>plot2d("f",0,1;0.4): // gambar dengan a=0.4
\end{eulerprompt}
\eulerimg{17}{images/EMT4Plot2D_Indonesia-029.png}
\begin{eulerprompt}
>plot2d(\{\{"f",0.2\}\},0,1): // gambar dengan a=0.2
\end{eulerprompt}
\eulerimg{17}{images/EMT4Plot2D_Indonesia-030.png}
\begin{eulerprompt}
>plot2d(\{\{"f(x,b)",b=0.1\}\},0,1): // gambar dengan 0.1
\end{eulerprompt}
\eulerimg{17}{images/EMT4Plot2D_Indonesia-031.png}
\begin{eulerprompt}
>function f(x) := x^3-x; ...
>plot2d("f",r=1):
\end{eulerprompt}
\eulerimg{17}{images/EMT4Plot2D_Indonesia-032.png}
\begin{eulercomment}
Berikut ringkasan fungsi yang diterima

- ekspresi atau ekspresi simbolik dalam x\\
- fungsi atau fungsi simbolik berdasarkan nama sebagai "f"\\
- fungsi simbolik cukup dengan namanya f

Fungsi plot2d() juga menerima fungsi simbolik. Untuk fungsi simbolik, cukup namanya
saja.
\end{eulercomment}
\begin{eulerprompt}
>function f(x) &= diff(x^x,x)
\end{eulerprompt}
\begin{euleroutput}
  
                              x
                             x  (log(x) + 1)
  
\end{euleroutput}
\begin{eulerprompt}
>plot2d(f,0,2):
\end{eulerprompt}
\eulerimg{17}{images/EMT4Plot2D_Indonesia-033.png}
\begin{eulercomment}
Tentu saja, untuk ekspresi atau ekspresi simbolik, nama variabelnya sudah cukup untuk
menggambarnya.
\end{eulercomment}
\begin{eulerprompt}
>expr &= sin(x)*exp(-x)
\end{eulerprompt}
\begin{euleroutput}
  
                                - x
                               E    sin(x)
  
\end{euleroutput}
\begin{eulerprompt}
>plot2d(expr,0,3pi):
\end{eulerprompt}
\eulerimg{17}{images/EMT4Plot2D_Indonesia-034.png}
\begin{eulerprompt}
>function f(x) &= x^x;
>plot2d(f,r=1,cx=1,cy=1,color=blue,thickness=2);
>plot2d(&diff(f(x),x),>add,color=red,style="-.-"):
\end{eulerprompt}
\eulerimg{17}{images/EMT4Plot2D_Indonesia-035.png}
\begin{eulercomment}
Untuk gaya garis tersedia berbagai opsi.

- style="...". Pilih dari "-", "--", "-.", ".", ".-.", "-.-".\\
- color: Lihat di bawah untuk warna.\\
- thickness: Defaultnya 1.

Warna dapat dipilih sebagai salah satu warna default, atau sebagai warna RGB.

- 0..15: indeks warna default.\\
- konstanta warna: white, black, red, green, blue, cyan, olive,
lightgray, gray, darkgray, orange, lightgreen, turquoise, lightblue,
lightorange, yellow\\
- rgb(red,green,blue): parameter berupa bilangan real pada [0,1].
\end{eulercomment}
\begin{eulerprompt}
>plot2d("exp(-x^2)",r=2,color=red,thickness=3,style="--"):
\end{eulerprompt}
\eulerimg{17}{images/EMT4Plot2D_Indonesia-036.png}
\begin{eulercomment}
Berikut tampilan warna-warna yang telah ditentukan di EMT.
\end{eulercomment}
\begin{eulerprompt}
>aspect(2); columnsplot(ones(1,16),lab=0:15,grid=0,color=0:15):
\end{eulerprompt}
\eulerimg{13}{images/EMT4Plot2D_Indonesia-037.png}
\begin{eulercomment}
Tetapi Anda dapat menggunakan warna apa pun.
\end{eulercomment}
\begin{eulerprompt}
>columnsplot(ones(1,16),grid=0,color=rgb(0,0,linspace(0,1,15))):
\end{eulerprompt}
\eulerimg{13}{images/EMT4Plot2D_Indonesia-038.png}
\eulerheading{Menggambar Beberapa Kurva pada bidang koordinat yang sama}
\begin{eulercomment}
Menggambar lebih dari satu fungsi (beberapa fungsi) ke dalam satu jendela dapat
dilakukan dengan berbagai cara. Salah satunya adalah menggunakan \textgreater{}add untuk beberapa
pemanggilan plot2d, kecuali pemanggilan pertama. Fitur ini sudah kita gunakan pada
contoh-contoh di atas.
\end{eulercomment}
\begin{eulerprompt}
>aspect(); plot2d("cos(x)",r=2,grid=6); plot2d("x",style=".",>add):
\end{eulerprompt}
\eulerimg{13}{images/EMT4Plot2D_Indonesia-039.png}
\begin{eulerprompt}
>aspect(1.5); plot2d("sin(x)",0,2pi); plot2d("cos(x)",color=blue,style="--",>add):
\end{eulerprompt}
\eulerimg{17}{images/EMT4Plot2D_Indonesia-040.png}
\begin{eulercomment}
Salah satu kegunaan \textgreater{}add adalah untuk menambahkan titik pada kurva.
\end{eulercomment}
\begin{eulerprompt}
>plot2d("sin(x)",0,pi); plot2d(2,sin(2),>points,>add):
\end{eulerprompt}
\eulerimg{17}{images/EMT4Plot2D_Indonesia-041.png}
\begin{eulercomment}
Kita tambahkan titik potong dengan label (pada posisi "cl" untuk center left, yaitu
tengah kiri), dan menyisipkan hasilnya ke notebook. Kita juga menambahkan judul pada
plot.
\end{eulercomment}
\begin{eulerprompt}
>plot2d(["cos(x)","x"],r=1.1,cx=0.5,cy=0.5, ...
>  color=[black,blue],style=["-","."], ...
>  grid=1);
>x0=solve("cos(x)-x",1);  ...
>  plot2d(x0,x0,>points,>add,title="Intersection Demo");  ...
>  label("cos(x) = x",x0,x0,pos="cl",offset=20):
\end{eulerprompt}
\eulerimg{17}{images/EMT4Plot2D_Indonesia-042.png}
\begin{eulercomment}
Pada demo berikut, kita menggambar fungsi sinc(x)=sin(x)/x beserta ekspansi Taylor
ke-8 dan ke-16. Ekspansi ini kita hitung menggunakan Maxima melalui ekspresi simbolik.
Plot ini dibuat dengan perintah multi-baris berikut dengan tiga pemanggilan plot2d().
Pemanggilan kedua dan ketiga memakai flag \textgreater{}add, yang membuat plot menggunakan rentang
sebelumnya.

Kita tambahkan kotak label yang menjelaskan fungsi-fungsi tersebut.
\end{eulercomment}
\begin{eulerprompt}
>$taylor(sin(x)/x,x,0,4)
\end{eulerprompt}
\begin{eulerformula}
\[
\frac{x^4}{120}-\frac{x^2}{6}+1
\]
\end{eulerformula}
\begin{eulerprompt}
>plot2d("sinc(x)",0,4pi,color=green,thickness=2); ...
>  plot2d(&taylor(sin(x)/x,x,0,8),>add,color=blue,style="--"); ...
>  plot2d(&taylor(sin(x)/x,x,0,16),>add,color=red,style="-.-"); ...
>  labelbox(["sinc","T8","T16"],styles=["-","--","-.-"], ...
>    colors=[black,blue,red]):
\end{eulerprompt}
\eulerimg{17}{images/EMT4Plot2D_Indonesia-044.png}
\begin{eulercomment}
Pada contoh berikut, kita membangkitkan Polinomial Bernstein.

\end{eulercomment}
\begin{eulerformula}
\[
B_i(x) = \binom{n}{i} x^i (1-x)^{n-i}
\]
\end{eulerformula}
\begin{eulerprompt}
>plot2d("(1-x)^10",0,1); // gambar fungsi pertama
>for i=1 to 10; plot2d("bin(10,i)*x^i*(1-x)^(10-i)",>add); end;
>insimg;
\end{eulerprompt}
\eulerimg{17}{images/EMT4Plot2D_Indonesia-045.png}
\begin{eulercomment}
Metode kedua adalah menggunakan pasangan matriks nilai-x dan matriks nilai-y yang
berukuran sama.

Kita membangkitkan matriks nilai dengan satu Polinomial Bernstein pada setiap
barisnya. Untuk itu, kita cukup menggunakan vektor kolom i. Lihat pengantar tentang
bahasa matriks untuk mempelajari rincian lebih lanjut.
\end{eulercomment}
\begin{eulerprompt}
>x=linspace(0,1,500);
>n=10; k=(0:n)'; // n adalah vektor baris, k adalah vektor kolom
>y=bin(n,k)*x^k*(1-x)^(n-k); // maka y adalah matriks
>plot2d(x,y):
\end{eulerprompt}
\eulerimg{17}{images/EMT4Plot2D_Indonesia-046.png}
\begin{eulercomment}
Perhatikan bahwa parameter color dapat berupa vektor. Dalam hal ini setiap warna
digunakan untuk setiap baris matriks.
\end{eulercomment}
\begin{eulerprompt}
>x=linspace(0,1,200); y=x^(1:10)'; plot2d(x,y,color=1:10):
\end{eulerprompt}
\eulerimg{17}{images/EMT4Plot2D_Indonesia-047.png}
\begin{eulercomment}
Metode lain adalah menggunakan vektor ekspresi (string). Anda kemudian dapat
menggunakan array warna, array gaya, dan array ketebalan dengan panjang yang sama.
\end{eulercomment}
\begin{eulerprompt}
>plot2d(["sin(x)","cos(x)"],0,2pi,color=4:5): 
\end{eulerprompt}
\eulerimg{17}{images/EMT4Plot2D_Indonesia-048.png}
\begin{eulerprompt}
>plot2d(["sin(x)","cos(x)"],0,2pi): // gambar vektor ekspresi
\end{eulerprompt}
\eulerimg{17}{images/EMT4Plot2D_Indonesia-049.png}
\begin{eulercomment}
Kita dapat memperoleh vektor seperti itu dari Maxima menggunakan makelist() dan
mxm2str().
\end{eulercomment}
\begin{eulerprompt}
>v &= makelist(binomial(10,i)*x^i*(1-x)^(10-i),i,0,10) // buat daftar
\end{eulerprompt}
\begin{euleroutput}
  
                 10            9              8  2             7  3
         [(1 - x)  , 10 (1 - x)  x, 45 (1 - x)  x , 120 (1 - x)  x , 
             6  4             5  5             4  6             3  7
  210 (1 - x)  x , 252 (1 - x)  x , 210 (1 - x)  x , 120 (1 - x)  x , 
            2  8              9   10
  45 (1 - x)  x , 10 (1 - x) x , x  ]
  
\end{euleroutput}
\begin{eulerprompt}
>mxm2str(v) // ambil vektor string dari vektor simbolik
\end{eulerprompt}
\begin{euleroutput}
  (1-x)^10
  10*(1-x)^9*x
  45*(1-x)^8*x^2
  120*(1-x)^7*x^3
  210*(1-x)^6*x^4
  252*(1-x)^5*x^5
  210*(1-x)^4*x^6
  120*(1-x)^3*x^7
  45*(1-x)^2*x^8
  10*(1-x)*x^9
  x^10
\end{euleroutput}
\begin{eulerprompt}
>plot2d(mxm2str(v),0,1): // gambar fungsi-fungsi
\end{eulerprompt}
\eulerimg{17}{images/EMT4Plot2D_Indonesia-050.png}
\begin{eulercomment}
Alternatif lain adalah menggunakan bahasa matriks Euler.

Jika sebuah ekspresi menghasilkan matriks fungsi, dengan satu fungsi pada setiap
barisnya, semua fungsi tersebut akan digambar dalam satu plot.

Untuk itu, gunakan vektor parameter berbentuk vektor kolom. Jika array warna
ditambahkan, array tersebut akan digunakan untuk setiap baris plot.
\end{eulercomment}
\begin{eulerprompt}
>n=(1:10)'; plot2d("x^n",0,1,color=1:10):
\end{eulerprompt}
\eulerimg{17}{images/EMT4Plot2D_Indonesia-051.png}
\begin{eulercomment}
Ekspresi dan fungsi satu baris dapat melihat variabel global.

Jika Anda tidak dapat menggunakan variabel global, Anda perlu menggunakan fungsi
dengan parameter tambahan, dan meneruskan parameter ini sebagai parameter titik koma.

Berhati-hatilah untuk menaruh semua parameter yang ditetapkan di bagian akhir perintah
plot2d. Pada contoh ini kita meneruskan a=5 ke fungsi f, yang kita gambar dari -10
sampai 10.
\end{eulercomment}
\begin{eulerprompt}
>function f(x,a) := 1/a*exp(-x^2/a); ...
>plot2d("f",-10,10;5,thickness=2,title="a=5"):
\end{eulerprompt}
\eulerimg{17}{images/EMT4Plot2D_Indonesia-052.png}
\begin{eulercomment}
Alternatifnya, gunakan koleksi yang berisi nama fungsi dan semua parameter tambahan.
Daftar khusus ini disebut call collection, dan itulah cara yang dianjurkan untuk
meneruskan argumen ke fungsi yang itu sendiri diteruskan sebagai argumen ke fungsi
lain.

Pada contoh berikut, kita menggunakan perulangan untuk menggambar beberapa fungsi
(lihat tutorial tentang pemrograman untuk perulangan).
\end{eulercomment}
\begin{eulerprompt}
>plot2d(\{\{"f",1\}\},-10,10); ...
>for a=2:10; plot2d(\{\{"f",a\}\},>add); end:
\end{eulerprompt}
\eulerimg{17}{images/EMT4Plot2D_Indonesia-053.png}
\begin{eulercomment}
Kita dapat mencapai hasil yang sama dengan cara berikut menggunakan bahasa matriks
EMT. Setiap baris matriks f(x,a) adalah satu fungsi. Selain itu, kita dapat mengatur
warna untuk setiap baris matriks. Klik dua kali pada fungsi getspectral() untuk
melihat penjelasannya.
\end{eulercomment}
\begin{eulerprompt}
>x=-10:0.01:10; a=(1:10)'; plot2d(x,f(x,a),color=getspectral(a/10)):
\end{eulerprompt}
\eulerimg{17}{images/EMT4Plot2D_Indonesia-054.png}
\eulersubheading{Label Teks}
\begin{eulercomment}
Dekorasi sederhana dapat berupa

- judul dengan title="..."\\
- label x dan y dengan xl="...", yl="..."\\
- label teks lain dengan label("...",x,y)

Perintah label akan menggambar pada plot saat ini di koordinat plot (x,y). Perintah
ini dapat menerima argumen posisional.
\end{eulercomment}
\begin{eulerprompt}
>plot2d("x^3-x",-1,2,title="y=x^3-x",yl="y",xl="x"):
\end{eulerprompt}
\eulerimg{17}{images/EMT4Plot2D_Indonesia-055.png}
\begin{eulerprompt}
>expr := "log(x)/x"; ...
>  plot2d(expr,0.5,5,title="y="+expr,xl="x",yl="y"); ...
>  label("(1,0)",1,0); label("Max",E,expr(E),pos="lc"):
\end{eulerprompt}
\eulerimg{17}{images/EMT4Plot2D_Indonesia-056.png}
\begin{eulercomment}
Ada juga fungsi labelbox(), yang dapat menampilkan fungsi-fungsi dan sebuah teks.
Fungsi ini menerima vektor string dan warna, satu item untuk setiap fungsi.
\end{eulercomment}
\begin{eulerprompt}
>function f(x) &= x^2*exp(-x^2);  ...
>plot2d(&f(x),a=-3,b=3,c=-1,d=1);  ...
>plot2d(&diff(f(x),x),>add,color=blue,style="--"); ...
>labelbox(["function","derivative"],styles=["-","--"], ...
>   colors=[black,blue],w=0.4):
\end{eulerprompt}
\eulerimg{17}{images/EMT4Plot2D_Indonesia-057.png}
\begin{eulercomment}
Kotak ini secara default ditambatkan di kanan atas, tetapi \textgreater{}left menambatkannya di
kiri atas. Anda dapat memindahkannya ke tempat mana pun yang Anda inginkan. Posisi
tambatan adalah sudut kanan atas kotak, dan angka-angkanya adalah pecahan dari ukuran
jendela grafik. Lebarnya ditentukan otomatis.

Untuk plot titik, kotak label juga berfungsi. Tambahkan parameter \textgreater{}points, atau vektor
flag, satu untuk setiap label.

Pada contoh berikut, hanya ada satu fungsi. Jadi kita dapat menggunakan string sebagai
ganti vektor string. Kita atur warna teks menjadi hitam untuk contoh ini.
\end{eulercomment}
\begin{eulerprompt}
>n=10; plot2d(0:n,bin(n,0:n),>addpoints); ...
>labelbox("Binomials",styles="[]",>points,x=0.1,y=0.1, ...
>tcolor=black,>left):
\end{eulerprompt}
\eulerimg{17}{images/EMT4Plot2D_Indonesia-058.png}
\begin{eulercomment}
Gaya plot ini juga tersedia di statplot(). Seperti pada plot2d(), warna dapat diatur
untuk setiap baris plot. Ada lebih banyak plot khusus untuk keperluan statistika
(lihat tutorial tentang statistika).
\end{eulercomment}
\begin{eulerprompt}
>statplot(1:10,random(2,10),color=[red,blue]):
\end{eulerprompt}
\eulerimg{17}{images/EMT4Plot2D_Indonesia-059.png}
\begin{eulercomment}
Fitur serupa adalah fungsi textbox().

Lebarnya secara default adalah lebar maksimal baris-baris teks. Tetapi lebar ini juga
dapat diatur oleh pengguna.
\end{eulercomment}
\begin{eulerprompt}
>function f(x) &= exp(-x)*sin(2*pi*x); ...
>plot2d("f(x)",0,2pi); ...
>textbox(latex("\(\backslash\)text\{Example of a damped oscillation\}\(\backslash\) f(x)=e^\{-x\}sin(2\(\backslash\)pi x)"),w=0.85):
\end{eulerprompt}
\eulerimg{17}{images/EMT4Plot2D_Indonesia-060.png}
\begin{eulercomment}
Label teks, judul, kotak label, dan teks lainnya dapat berisi string Unicode (lihat
sintaks EMT untuk lebih lanjut tentang string Unicode).
\end{eulercomment}
\begin{eulerprompt}
>plot2d("x^3-x",title=u"x &rarr; x&sup3; - x"):
\end{eulerprompt}
\eulerimg{17}{images/EMT4Plot2D_Indonesia-061.png}
\begin{eulercomment}
Label pada sumbu x dan y dapat vertikal, demikian pula sumbunya.
\end{eulercomment}
\begin{eulerprompt}
>plot2d("sinc(x)",0,2pi,xl="x",yl=u"x &rarr; sinc(x)",>vertical):
\end{eulerprompt}
\eulerimg{17}{images/EMT4Plot2D_Indonesia-062.png}
\eulersubheading{LaTeX}
\begin{eulercomment}
Anda juga dapat menggambar rumus LaTeX jika Anda telah memasang sistem LaTeX. Saya
merekomendasikan MiKTeX. Path ke program "latex" dan "dvipng" harus ada pada path
sistem, atau Anda harus mengatur LaTeX pada menu opsi.

Perhatikan bahwa parsing LaTeX lambat. Jika Anda ingin menggunakan LaTeX pada plot
animasi, panggillah latex() satu kali sebelum perulangan dan gunakan hasilnya (sebuah
gambar dalam matriks RGB).

Pada plot berikut, kita menggunakan LaTeX untuk label x dan y, sebuah label, sebuah
kotak label, dan judul plot.
\end{eulercomment}
\begin{eulerprompt}
>plot2d("exp(-x)*sin(x)/x",a=0,b=2pi,c=0,d=1,grid=6,color=blue, ...
>  title=latex("\(\backslash\)text\{Function $\(\backslash\)Phi$\}"), ...
>  xl=latex("\(\backslash\)phi"),yl=latex("\(\backslash\)Phi(\(\backslash\)phi)")); ...
>textbox( ...
>  latex("\(\backslash\)Phi(\(\backslash\)phi) = e^\{-\(\backslash\)phi\} \(\backslash\)frac\{\(\backslash\)sin(\(\backslash\)phi)\}\{\(\backslash\)phi\}"),x=0.8,y=0.5); ...
>label(latex("\(\backslash\)Phi",color=blue),1,0.4):
\end{eulerprompt}
\eulerimg{17}{images/EMT4Plot2D_Indonesia-063.png}
\begin{eulercomment}
Sering kali kita menginginkan jarak yang tidak konformal dan label teks pada sumbu x.
Kita dapat menggunakan xaxis() dan yaxis() seperti yang akan ditunjukkan nanti.

Cara termudah adalah membuat plot kosong dengan bingkai menggunakan grid=4, lalu
menambahkan grid dengan ygrid() dan xgrid(). Pada contoh berikut, kita menggunakan
tiga string LaTeX untuk label pada sumbu x dengan xtick().
\end{eulercomment}
\begin{eulerprompt}
>plot2d("sinc(x)",0,2pi,grid=4,<ticks); ...
>ygrid(-2:0.5:2,grid=6); ...
>xgrid([0:2]*pi,<ticks,grid=6);  ...
>xtick([0,pi,2pi],["0","\(\backslash\)pi","2\(\backslash\)pi"],>latex):
\end{eulerprompt}
\eulerimg{17}{images/EMT4Plot2D_Indonesia-064.png}
\begin{eulercomment}
Tentu saja, fungsi juga dapat digunakan.
\end{eulercomment}
\begin{eulerprompt}
>function map f(x) ...
\end{eulerprompt}
\begin{eulerudf}
  if x>0 then return x^4
  else return x^2
  endif
  endfunction
\end{eulerudf}
\begin{eulercomment}
Parameter "map" membantu menggunakan fungsi untuk vektor. Untuk plotnya sendiri, hal
ini tidak diperlukan. Tetapi untuk menunjukkan bahwa vektorisasi berguna, kita
menambahkan beberapa titik kunci pada plot di x=-1, x=0 dan x=1.

Pada plot berikut, kita juga memasukkan beberapa kode LaTeX. Kita menggunakannya untuk
dua label dan sebuah kotak teks. Tentu saja, Anda hanya dapat menggunakan LaTeX jika
LaTeX telah terpasang dengan benar.
\end{eulercomment}
\begin{eulerprompt}
>plot2d("f",-1,1,xl="x",yl="f(x)",grid=6);  ...
>plot2d([-1,0,1],f([-1,0,1]),>points,>add); ...
>label(latex("x^3"),0.72,f(0.72)); ...
>label(latex("x^2"),-0.52,f(-0.52),pos="ll"); ...
>textbox( ...
>  latex("f(x)=\(\backslash\)begin\{cases\} x^3 & x>0 \(\backslash\)\(\backslash\) x^2 & x \(\backslash\)le 0\(\backslash\)end\{cases\}"), ...
>  x=0.7,y=0.2):
\end{eulerprompt}
\eulerimg{17}{images/EMT4Plot2D_Indonesia-065.png}
\begin{eulercomment}
\end{eulercomment}
\eulersubheading{Interaksi Pengguna}
\begin{eulercomment}
Ketika menggambar fungsi atau ekspresi, parameter \textgreater{}user memungkinkan pengguna
memperbesar dan menggeser plot dengan tombol kursor atau mouse. Pengguna dapat

- memperbesar/memperkecil dengan + atau -\\
- menggeser plot dengan tombol kursor\\
- memilih jendela plot dengan mouse\\
- mengatur ulang tampilan dengan spasi\\
- keluar dengan return

Tombol spasi akan mengatur ulang plot ke jendela plot semula.

Ketika menggambar data, flag \textgreater{}user hanya akan menunggu penekanan tombol.
\end{eulercomment}
\begin{eulerprompt}
>plot2d(\{\{"x^3-a*x",a=1\}\},>user,title="Press any key!"):
\end{eulerprompt}
\eulerimg{17}{images/EMT4Plot2D_Indonesia-066.png}
\begin{eulerprompt}
>plot2d("exp(x)*sin(x)",user=true, ...
>  title="+/- or cursor keys (return to exit)"):
\end{eulerprompt}
\eulerimg{17}{images/EMT4Plot2D_Indonesia-067.png}
\begin{eulercomment}
Berikut ini mendemonstrasikan cara interaksi pengguna yang lebih lanjut (lihat
tutorial tentang pemrograman untuk rincian).

Fungsi bawaan mousedrag() menunggu peristiwa mouse atau keyboard. Fungsi ini
melaporkan mouse ditekan, mouse digerakkan atau mouse dilepas, serta penekanan tombol.
Fungsi dragpoints() memanfaatkan hal ini, dan memungkinkan pengguna menyeret titik
mana pun pada plot.

Pertama-tama kita memerlukan fungsi plot. Sebagai contoh, kita melakukan interpolasi
pada 5 titik dengan polinomial. Fungsi tersebut harus menggambar pada area plot yang
tetap.
\end{eulercomment}
\begin{eulerprompt}
>function plotf(xp,yp,select) ...
\end{eulerprompt}
\begin{eulerudf}
    d=interp(xp,yp);
    plot2d("interpval(xp,d,x)";d,xp,r=2);
    plot2d(xp,yp,>points,>add);
    if select>0 then
      plot2d(xp[select],yp[select],color=red,>points,>add);
    endif;
    title("Drag one point, or press space or return!");
  endfunction
\end{eulerudf}
\begin{eulercomment}
Perhatikan parameter titik koma pada plot2d (d dan xp), yang diteruskan ke evaluasi
fungsi interp(). Tanpa ini, kita harus menulis fungsi plotinterp() terlebih dahulu,
yang mengakses nilai-nilai secara global.

Sekarang kita membangkitkan beberapa nilai acak, dan membiarkan pengguna menyeret
titik-titiknya.
\end{eulercomment}
\begin{eulerprompt}
>t=-1:0.5:1; dragpoints("plotf",t,random(size(t))-0.5):
\end{eulerprompt}
\eulerimg{17}{images/EMT4Plot2D_Indonesia-068.png}
\begin{eulercomment}
Ada juga fungsi yang menggambar fungsi lain bergantung pada vektor parameter, dan
membiarkan pengguna menyesuaikan parameter-parameter tersebut.

Pertama kita memerlukan fungsi plotnya.
\end{eulercomment}
\begin{eulerprompt}
>function plotf([a,b]) := plot2d("exp(a*x)*cos(2pi*b*x)",0,2pi;a,b);
\end{eulerprompt}
\begin{eulercomment}
Kemudian kita memerlukan nama-nama parameter, nilai awal, dan matriks rentang nx2,
secara opsional sebuah baris judul. Terdapat slider interaktif, yang dapat mengatur
nilai oleh pengguna. Fungsi dragvalues() menyediakannya.
\end{eulercomment}
\begin{eulerprompt}
>dragvalues("plotf",["a","b"],[-1,2],[[-2,2];[1,10]], ...
>  heading="Drag these values:",hcolor=black):
\end{eulerprompt}
\eulerimg{17}{images/EMT4Plot2D_Indonesia-069.png}
\begin{eulercomment}
Nilai yang diseret dapat dibatasi pada bilangan bulat. Sebagai contoh, kita menulis
fungsi plot yang menggambar polinomial Taylor berderajat n untuk fungsi kosinus.
\end{eulercomment}
\begin{eulerprompt}
>function plotf(n) ...
\end{eulerprompt}
\begin{eulerudf}
  plot2d("cos(x)",0,2pi,>square,grid=6);
  plot2d(&"taylor(cos(x),x,0,@n)",color=blue,>add);
  textbox("Taylor polynomial of degree "+n,0.1,0.02,style="t",>left);
  endfunction
\end{eulerudf}
\begin{eulercomment}
Sekarang kita biarkan derajat n bervariasi dari 0 sampai 20 dalam 20 langkah. Hasil
dragvalues() digunakan untuk menggambar sketsa dengan n tersebut, dan menyisipkan plot
ke notebook.
\end{eulercomment}
\begin{eulerprompt}
>nd=dragvalues("plotf","degree",2,[0,20],20,y=0.8, ...
>   heading="Drag the value:"); ...
>plotf(nd):
\end{eulerprompt}
\eulerimg{17}{images/EMT4Plot2D_Indonesia-070.png}
\begin{eulercomment}
Berikut adalah demonstrasi sederhana dari fungsi ini. Pengguna dapat menggambar di
atas jendela plot, meninggalkan jejak titik-titik.
\end{eulercomment}
\begin{eulerprompt}
>function dragtest ...
\end{eulerprompt}
\begin{eulerudf}
    plot2d(none,r=1,title="Drag with the mouse, or press any key!");
    start=0;
    repeat
      \{flag,m,time\}=mousedrag();
      if flag==0 then return; endif;
      if flag==2 then
        hold on; mark(m[1],m[2]); hold off;
      endif;
    end
  endfunction
\end{eulerudf}
\begin{eulerprompt}
>dragtest // lihat hasilnya dan cobalah lakukan!
\end{eulerprompt}
\eulersubheading{Gaya Plot 2D}
\begin{eulercomment}
Secara default, EMT menghitung tanda sumbu (tick) secara otomatis dan menambahkan
label pada setiap tick. Hal ini dapat diubah dengan parameter grid. Gaya default sumbu
dan label dapat dimodifikasi. Selain itu, label dan judul dapat ditambahkan secara
manual. Untuk mengembalikan ke gaya default, gunakan reset().
\end{eulercomment}
\begin{eulerprompt}
>aspect();
>figure(3,4); ...
> figure(1); plot2d("x^3-x",grid=0); ... // tanpa grid, bingkai, atau sumbu
> figure(2); plot2d("x^3-x",grid=1); ... // sumbu x-y
> figure(3); plot2d("x^3-x",grid=2); ... // tick default
> figure(4); plot2d("x^3-x",grid=3); ... // sumbu x-y dengan label di dalam
> figure(5); plot2d("x^3-x",grid=4); ... // tanpa tick, hanya label
> figure(6); plot2d("x^3-x",grid=5); ... // default, tetapi tanpa margin
> figure(7); plot2d("x^3-x",grid=6); ... // hanya sumbu
> figure(8); plot2d("x^3-x",grid=7); ... // hanya sumbu, tick pada sumbu
> figure(9); plot2d("x^3-x",grid=8); ... // hanya sumbu, tick lebih halus pada sumbu
> figure(10); plot2d("x^3-x",grid=9); ... // default, tick kecil di dalam
> figure(11); plot2d("x^3-x",grid=10); ...// no ticks, axes only
> figure(0):
\end{eulerprompt}
\eulerimg{13}{images/EMT4Plot2D_Indonesia-071.png}
\begin{eulercomment}
Parameter \textless{}frame mematikan bingkai, dan framecolor=blue mengatur bingkai menjadi
berwarna biru.

Jika Anda menginginkan tick sendiri, Anda dapat menggunakan style=0, dan menambahkan
semuanya kemudian.
\end{eulercomment}
\begin{eulerprompt}
>aspect(1.5); 
>plot2d("x^3-x",grid=0); // gambar
>frame; xgrid([-1,0,1]); ygrid(0): // tambahkan bingkai dan grid
\end{eulerprompt}
\eulerimg{17}{images/EMT4Plot2D_Indonesia-072.png}
\begin{eulercomment}
Untuk judul plot dan label sumbu, lihat contoh berikut.
\end{eulercomment}
\begin{eulerprompt}
>plot2d("exp(x)",-1,1);
>textcolor(black); // atur warna teks menjadi hitam
>title(latex("y=e^x")); // judul di atas plot
>xlabel(latex("x")); // "x" untuk sumbu x
>ylabel(latex("y"),>vertical); // "y" vertikal untuk sumbu y
>label(latex("(0,1)"),0,1,color=blue): // beri label pada sebuah titik
\end{eulerprompt}
\eulerimg{17}{images/EMT4Plot2D_Indonesia-073.png}
\begin{eulercomment}
Sumbu dapat digambar secara terpisah dengan xaxis() dan yaxis().
\end{eulercomment}
\begin{eulerprompt}
>plot2d("x^3-x",<grid,<frame);
>xaxis(0,xx=-2:1,style="->"); yaxis(0,yy=-5:5,style="->"):
\end{eulerprompt}
\eulerimg{17}{images/EMT4Plot2D_Indonesia-074.png}
\begin{eulercomment}
Teks pada plot dapat diatur dengan label(). Pada contoh berikut, "lc" berarti lower
center (tengah bawah). Ini mengatur posisi label relatif terhadap koordinat plot.
\end{eulercomment}
\begin{eulerprompt}
>function f(x) &= x^3-x
\end{eulerprompt}
\begin{euleroutput}
  
                                   3
                                  x  - x
  
\end{euleroutput}
\begin{eulerprompt}
>plot2d(f,-1,1,>square);
>x0=fmin(f,0,1); // hitung titik minimum
>label("Rel. Min.",x0,f(x0),pos="lc"): // tambahkan label di sana
\end{eulerprompt}
\eulerimg{17}{images/EMT4Plot2D_Indonesia-075.png}
\begin{eulercomment}
Ada juga kotak teks.
\end{eulercomment}
\begin{eulerprompt}
>plot2d(&f(x),-1,1,-2,2); // fungsi
>plot2d(&diff(f(x),x),>add,style="--",color=red); // turunan
>labelbox(["f","f'"],["-","--"],[black,red]): // kotak label
\end{eulerprompt}
\eulerimg{17}{images/EMT4Plot2D_Indonesia-076.png}
\begin{eulerprompt}
>plot2d(["exp(x)","1+x"],color=[black,blue],style=["-","-.-"]):
\end{eulerprompt}
\eulerimg{17}{images/EMT4Plot2D_Indonesia-077.png}
\begin{eulerprompt}
>gridstyle("->",color=gray,textcolor=gray,framecolor=gray);  ...
> plot2d("x^3-x",grid=1);   ...
> settitle("y=x^3-x",color=black); ...
> label("x",2,0,pos="bc",color=gray);  ...
> label("y",0,6,pos="cl",color=gray); ...
> reset():
\end{eulerprompt}
\eulerimg{13}{images/EMT4Plot2D_Indonesia-078.png}
\begin{eulercomment}
Untuk kontrol yang lebih besar lagi, sumbu x dan sumbu y dapat dibuat secara manual.

Perintah fullwindow() memperluas jendela plot karena kita tidak lagi memerlukan ruang
untuk label di luar jendela plot. Gunakan shrinkwindow() atau reset() untuk
mengembalikan ke pengaturan default.
\end{eulercomment}
\begin{eulerprompt}
>fullwindow; ...
> gridstyle(color=darkgray,textcolor=darkgray); ...
> plot2d(["2^x","1","2^(-x)"],a=-2,b=2,c=0,d=4,<grid,color=4:6,<frame); ...
> xaxis(0,-2:1,style="->"); xaxis(0,2,"x",<axis); ...
> yaxis(0,4,"y",style="->"); ...
> yaxis(-2,1:4,>left); ...
> yaxis(2,2^(-2:2),style=".",<left); ...
> labelbox(["2^x","1","2^-x"],colors=4:6,x=0.8,y=0.2); ...
> reset:
\end{eulerprompt}
\eulerimg{13}{images/EMT4Plot2D_Indonesia-079.png}
\begin{eulercomment}
Berikut contoh lain, di mana string Unicode digunakan dan sumbu berada di luar area
plot.
\end{eulercomment}
\begin{eulerprompt}
>aspect(1.5); 
>plot2d(["sin(x)","cos(x)"],0,2pi,color=[red,green],<grid,<frame); ...
> xaxis(-1.1,(0:2)*pi,xt=["0",u"&pi;",u"2&pi;"],style="-",>ticks,>zero);  ...
> xgrid((0:0.5:2)*pi,<ticks); ...
> yaxis(-0.1*pi,-1:0.2:1,style="-",>zero,>grid); ...
> labelbox(["sin","cos"],colors=[red,green],x=0.5,y=0.2,>left); ...
> xlabel(u"&phi;"); ylabel(u"f(&phi;)"):
\end{eulerprompt}
\eulerimg{17}{images/EMT4Plot2D_Indonesia-080.png}
\eulerheading{Menggambar Data 2D}
\begin{eulercomment}
Jika x dan y adalah vektor data, data ini akan digunakan sebagai koordinat x dan y
suatu kurva. Dalam hal ini, a, b, c, dan d, atau jari-jari r dapat ditentukan, atau
jendela plot akan menyesuaikan secara otomatis dengan data. Alternatifnya, \textgreater{}square
dapat diatur untuk mempertahankan rasio aspek persegi.

Menggambar ekspresi hanyalah singkatan untuk plot data. Untuk plot data, Anda
memerlukan satu atau lebih baris nilai x, dan satu atau lebih baris nilai y. Dari
rentang dan nilai-nilai x, fungsi plot2d akan menghitung data yang digambar, secara
default dengan evaluasi fungsi adaptif. Untuk plot titik gunakan "\textgreater{}points", untuk
campuran garis dan titik gunakan "\textgreater{}addpoints".

Tetapi Anda dapat memasukkan data secara langsung.

- Gunakan vektor baris untuk x dan y untuk satu fungsi.\\
- Matriks untuk x dan y digambar baris demi baris.

Berikut contoh dengan satu baris untuk x dan y.

\end{eulercomment}
\begin{eulerprompt}
>x=-10:0.1:10; y=exp(-x^2)*x; plot2d(x,y):
\end{eulerprompt}
\eulerimg{17}{images/EMT4Plot2D_Indonesia-081.png}
\begin{eulercomment}
Data juga dapat digambar sebagai titik. Gunakan points=true untuk ini. Plot bekerja
seperti poligon, tetapi hanya menggambar sudut-sudutnya.

- style="...": Pilih dari "[]", "\textless{}\textgreater{}", "o", ".", "..", "+", "*", "[]#", "\textless{}\textgreater{}#", "o#",
"..#", "#", "\textbar{}".

Untuk menggambar himpunan titik gunakan \textgreater{}points. Jika warna berupa vektor warna,
setiap titik mendapat warna berbeda. Untuk matriks koordinat dan vektor kolom, warna
berlaku untuk baris-baris matriks. Parameter \textgreater{}addpoints menambahkan titik pada ruas
garis untuk plot data.
\end{eulercomment}
\begin{eulerprompt}
>xdata=[1,1.5,2.5,3,4]; ydata=[3,3.1,2.8,2.9,2.7]; // data
>plot2d(xdata,ydata,a=0.5,b=4.5,c=2.5,d=3.5,style="."); // garis
>plot2d(xdata,ydata,>points,>add,style="o"): // tambahkan titik
\end{eulerprompt}
\eulerimg{17}{images/EMT4Plot2D_Indonesia-082.png}
\begin{eulerprompt}
>p=polyfit(xdata,ydata,1); // hitung garis regresi
>plot2d("polyval(p,x)",>add,color=red): // tambahkan plot garis
\end{eulerprompt}
\eulerimg{17}{images/EMT4Plot2D_Indonesia-083.png}
\eulerheading{Menggambar Daerah Yang Dibatasi Kurva}
\begin{eulercomment}
Plot data sebenarnya adalah poligon. Kita juga dapat menggambar kurva atau kurva
berarsir (terisi).

- filled=true mengisi plot.\\
- style="...": Pilih dari "#", "/", "\textbackslash{}", "\textbackslash{}/".\\
- fillcolor: Lihat di atas untuk warna yang tersedia.

Warna isian ditentukan oleh argumen "fillcolor", dan \textless{}outline opsional mencegah
penggambaran batas untuk semua gaya kecuali gaya default.
\end{eulercomment}
\begin{eulerprompt}
>t=linspace(0,2pi,1000); // parameter untuk kurva
>x=sin(t)*exp(t/pi); y=cos(t)*exp(t/pi); // x(t) dan y(t)
>figure(1,2); aspect(16/9)
>figure(1); plot2d(x,y,r=10); // gambar kurva
>figure(2); plot2d(x,y,r=10,>filled,style="/",fillcolor=red); // isi kurva
>figure(0):
\end{eulerprompt}
\eulerimg{14}{images/EMT4Plot2D_Indonesia-084.png}
\begin{eulercomment}
Pada contoh berikut kita menggambar elips terisi dan dua segi enam terisi menggunakan
kurva tertutup dengan 6 titik dengan gaya isian yang berbeda.
\end{eulercomment}
\begin{eulerprompt}
>x=linspace(0,2pi,1000); plot2d(sin(x),cos(x)*0.5,r=1,>filled,style="/"):
\end{eulerprompt}
\eulerimg{14}{images/EMT4Plot2D_Indonesia-085.png}
\begin{eulerprompt}
>t=linspace(0,2pi,6); ...
>plot2d(cos(t),sin(t),>filled,style="/",fillcolor=red,r=1.2):
\end{eulerprompt}
\eulerimg{14}{images/EMT4Plot2D_Indonesia-086.png}
\begin{eulerprompt}
>t=linspace(0,2pi,6); plot2d(cos(t),sin(t),>filled,style="#"):
\end{eulerprompt}
\eulerimg{14}{images/EMT4Plot2D_Indonesia-087.png}
\begin{eulercomment}
Contoh lain adalah segi tujuh, yang kita buat dengan 7 titik pada lingkaran satuan.
\end{eulercomment}
\begin{eulerprompt}
>t=linspace(0,2pi,7);  ...
> plot2d(cos(t),sin(t),r=1,>filled,style="/",fillcolor=red):
\end{eulerprompt}
\eulerimg{14}{images/EMT4Plot2D_Indonesia-088.png}
\begin{eulercomment}
Berikut adalah himpunan nilai maksimal dari empat kondisi linear yang kurang dari atau
sama dengan 3. Ini adalah A[k].v\textless{}=3 untuk semua baris A. Untuk mendapatkan sudut yang
bagus, kita menggunakan n yang relatif besar.
\end{eulercomment}
\begin{eulerprompt}
>A=[2,1;1,2;-1,0;0,-1];
>function f(x,y) := max([x,y].A');
>plot2d("f",r=4,level=[0;3],color=green,n=111):
\end{eulerprompt}
\eulerimg{14}{images/EMT4Plot2D_Indonesia-089.png}
\begin{eulercomment}
Inti dari bahasa matriks adalah bahwa bahasa ini memudahkan pembuatan tabel fungsi.
\end{eulercomment}
\begin{eulerprompt}
>t=linspace(0,2pi,1000); x=cos(3*t); y=sin(4*t);
\end{eulerprompt}
\begin{eulercomment}
Sekarang kita memiliki vektor x dan y berisi nilai-nilai. plot2d() dapat menggambar\\
nilai-nilai ini sebagai kurva yang menghubungkan titik-titik. Plot dapat diisi. Dalam\\
hal ini hasilnya bagus karena aturan lilitan (winding rule) yang digunakan untuk\\
pengisian.
\end{eulercomment}
\begin{eulerprompt}
>plot2d(x,y,<grid,<frame,>filled):
\end{eulerprompt}
\eulerimg{14}{images/EMT4Plot2D_Indonesia-090.png}
\begin{eulercomment}
Vektor interval digambar terhadap nilai x sebagai daerah terisi di antara\\
nilai bawah dan atas interval.

Ini dapat berguna untuk menggambar galat perhitungan. Tetapi dapat juga\\
digunakan untuk menggambar galat statistik.
\end{eulercomment}
\begin{eulerprompt}
>t=0:0.1:1; ...
> plot2d(t,interval(t-random(size(t)),t+random(size(t))),style="|");  ...
> plot2d(t,t,add=true):
\end{eulerprompt}
\eulerimg{14}{images/EMT4Plot2D_Indonesia-091.png}
\begin{eulercomment}
Jika x adalah vektor terurut, dan y adalah vektor interval, maka plot2d akan
menggambar rentang interval yang terisi pada bidang. Gaya isian sama dengan gaya
poligon.
\end{eulercomment}
\begin{eulerprompt}
>t=-1:0.01:1; x=~t-0.01,t+0.01~; y=x^3-x;
>plot2d(t,y):
\end{eulerprompt}
\eulerimg{14}{images/EMT4Plot2D_Indonesia-092.png}
\begin{eulercomment}
Dimungkinkan untuk mengisi daerah nilai untuk fungsi tertentu. Untuk\\
ini, level harus berupa matriks 2xn. Baris pertama adalah batas bawah\\
dan baris kedua berisi batas atas.
\end{eulercomment}
\begin{eulerprompt}
>expr := "2*x^2+x*y+3*y^4+y"; // definisikan ekspresi f(x,y)
>plot2d(expr,level=[0;1],style="-",color=blue): // 0 <= f(x,y) <= 1
\end{eulerprompt}
\eulerimg{14}{images/EMT4Plot2D_Indonesia-093.png}
\begin{eulercomment}
Kita juga dapat mengisi rentang nilai seperti

\end{eulercomment}
\begin{eulerformula}
\[
-1 \le (x^2+y^2)^2-x^2+y^2 \le 0.
\]
\end{eulerformula}
\begin{eulercomment}
\end{eulercomment}
\begin{eulerprompt}
>plot2d("(x^2+y^2)^2-x^2+y^2",r=1.2,level=[-1;0],style="/"):
\end{eulerprompt}
\eulerimg{14}{images/EMT4Plot2D_Indonesia-094.png}
\begin{eulerprompt}
>plot2d("cos(x)","sin(x)^3",xmin=0,xmax=2pi,>filled,style="/"):
\end{eulerprompt}
\eulerimg{14}{images/EMT4Plot2D_Indonesia-095.png}
\eulerheading{Grafik Fungsi Parametrik}
\begin{eulercomment}
Nilai-nilai x tidak perlu terurut. (x,y) hanya menggambarkan sebuah kurva. Jika x
terurut, kurva tersebut adalah grafik suatu fungsi.

Pada contoh berikut, kita menggambar spiral

\end{eulercomment}
\begin{eulerformula}
\[
\gamma(t) = t \cdot (\cos(2\pi t),\sin(2\pi t))
\]
\end{eulerformula}
\begin{eulercomment}
Kita perlu menggunakan sangat banyak titik agar tampak mulus, atau menggunakan fungsi
adaptive() untuk mengevaluasi ekspresi (lihat fungsi adaptive() untuk rincian lebih
lanjut).
\end{eulercomment}
\begin{eulerprompt}
>t=linspace(0,1,1000); ...
>plot2d(t*cos(2*pi*t),t*sin(2*pi*t),r=1):
\end{eulerprompt}
\eulerimg{14}{images/EMT4Plot2D_Indonesia-096.png}
\begin{eulercomment}
Alternatifnya, dimungkinkan menggunakan dua ekspresi untuk kurva. Berikut ini
menggambar kurva yang sama seperti di atas.
\end{eulercomment}
\begin{eulerprompt}
>plot2d("x*cos(2*pi*x)","x*sin(2*pi*x)",xmin=0,xmax=1,r=1):
\end{eulerprompt}
\eulerimg{14}{images/EMT4Plot2D_Indonesia-097.png}
\begin{eulerprompt}
>t=linspace(0,1,1000); r=exp(-t); x=r*cos(2pi*t); y=r*sin(2pi*t);
>plot2d(x,y,r=1):
\end{eulerprompt}
\eulerimg{14}{images/EMT4Plot2D_Indonesia-098.png}
\begin{eulercomment}
Pada contoh berikut, kita menggambar kurva

\end{eulercomment}
\begin{eulerformula}
\[
\gamma(t) = (r(t) \cos(t), r(t) \sin(t))
\]
\end{eulerformula}
\begin{eulercomment}
dengan

\end{eulercomment}
\begin{eulerformula}
\[
r(t) = 1 + \dfrac{\sin(3t)}{2}.
\]
\end{eulerformula}
\begin{eulerprompt}
>t=linspace(0,2pi,1000); r=1+sin(3*t)/2; x=r*cos(t); y=r*sin(t); ...
>plot2d(x,y,>filled,fillcolor=red,style="/",r=1.5):
\end{eulerprompt}
\eulerimg{14}{images/EMT4Plot2D_Indonesia-099.png}
\eulerheading{Menggambar Grafik Bilangan Kompleks}
\begin{eulercomment}
Array bilangan kompleks juga dapat digambar. Dalam hal ini titik-titik grid akan
dihubungkan. Jika banyaknya garis grid ditentukan (atau vektor 1x2 garis grid) pada
argumen cgrid, hanya garis-garis grid tersebut yang terlihat.

Matriks bilangan kompleks akan otomatis digambar sebagai grid pada bidang kompleks.

Pada contoh berikut, kita menggambar bayangan lingkaran satuan di bawah fungsi
eksponensial. Parameter cgrid menyembunyikan sebagian kurva grid.
\end{eulercomment}
\begin{eulerprompt}
>aspect(); r=linspace(0,1,50); a=linspace(0,2pi,80)'; z=r*exp(I*a);...
>plot2d(z,a=-1.25,b=1.25,c=-1.25,d=1.25,cgrid=10):
\end{eulerprompt}
\eulerimg{13}{images/EMT4Plot2D_Indonesia-100.png}
\begin{eulerprompt}
>aspect(1.25); r=linspace(0,1,50); a=linspace(0,2pi,200)'; z=r*exp(I*a);
>plot2d(exp(z),cgrid=[40,10]):
\end{eulerprompt}
\eulerimg{21}{images/EMT4Plot2D_Indonesia-101.png}
\begin{eulerprompt}
>r=linspace(0,1,10); a=linspace(0,2pi,40)'; z=r*exp(I*a);
>plot2d(exp(z),>points,>add):
\end{eulerprompt}
\eulerimg{21}{images/EMT4Plot2D_Indonesia-102.png}
\begin{eulercomment}
Vektor bilangan kompleks secara otomatis digambar sebagai kurva pada bidang kompleks
dengan bagian real dan bagian imajiner.

Pada contoh, kita menggambar lingkaran satuan dengan

\end{eulercomment}
\begin{eulerformula}
\[
\gamma(t) = e^{it}
\]
\end{eulerformula}
\begin{eulerprompt}
>t=linspace(0,2pi,1000); ...
>plot2d(exp(I*t)+exp(4*I*t),r=2):
\end{eulerprompt}
\eulerimg{21}{images/EMT4Plot2D_Indonesia-103.png}
\eulerheading{Plot Statistik}
\begin{eulercomment}
Ada banyak fungsi yang dikhususkan untuk plot statistik. Salah satu plot yang sering
digunakan adalah plot kolom.

Jumlah kumulatif dari nilai-nilai berdistribusi normal 0-1 menghasilkan jalan acak
(random walk).
\end{eulercomment}
\begin{eulerprompt}
>plot2d(cumsum(randnormal(1,1000))):
\end{eulerprompt}
\eulerimg{21}{images/EMT4Plot2D_Indonesia-104.png}
\begin{eulercomment}
Menggunakan dua baris menampilkan jalan acak dalam dua dimensi.
\end{eulercomment}
\begin{eulerprompt}
>X=cumsum(randnormal(2,1000)); plot2d(X[1],X[2]):
\end{eulerprompt}
\eulerimg{21}{images/EMT4Plot2D_Indonesia-105.png}
\begin{eulerprompt}
>columnsplot(cumsum(random(10)),style="/",color=blue):
\end{eulerprompt}
\eulerimg{21}{images/EMT4Plot2D_Indonesia-106.png}
\begin{eulercomment}
Plot ini juga dapat menampilkan string sebagai label.
\end{eulercomment}
\begin{eulerprompt}
>months=["Jan","Feb","Mar","Apr","May","Jun", ...
>  "Jul","Aug","Sep","Oct","Nov","Dec"];
>values=[10,12,12,18,22,28,30,26,22,18,12,8];
>columnsplot(values,lab=months,color=red,style="-");
>title("Temperature"):
\end{eulerprompt}
\eulerimg{21}{images/EMT4Plot2D_Indonesia-107.png}
\begin{eulerprompt}
>k=0:10;
>plot2d(k,bin(10,k),>bar):
\end{eulerprompt}
\eulerimg{21}{images/EMT4Plot2D_Indonesia-108.png}
\begin{eulerprompt}
>plot2d(k,bin(10,k)); plot2d(k,bin(10,k),>points,>add):
\end{eulerprompt}
\eulerimg{21}{images/EMT4Plot2D_Indonesia-109.png}
\begin{eulerprompt}
>plot2d(normal(1000),normal(1000),>points,grid=6,style=".."):
\end{eulerprompt}
\eulerimg{21}{images/EMT4Plot2D_Indonesia-110.png}
\begin{eulerprompt}
>plot2d(normal(1,1000),>distribution,style="O"):
\end{eulerprompt}
\eulerimg{21}{images/EMT4Plot2D_Indonesia-111.png}
\begin{eulerprompt}
>plot2d("qnormal",0,5;2.5,0.5,>filled):
\end{eulerprompt}
\eulerimg{21}{images/EMT4Plot2D_Indonesia-112.png}
\begin{eulercomment}
Untuk menggambar distribusi statistik eksperimental, Anda dapat menggunakan
distribution=n pada plot2d.
\end{eulercomment}
\begin{eulerprompt}
>w=randexponential(1,1000); // distribusi eksponensial
>plot2d(w,>distribution): // atau distribution=n dengan n interval
\end{eulerprompt}
\eulerimg{21}{images/EMT4Plot2D_Indonesia-113.png}
\begin{eulercomment}
Atau Anda dapat menghitung distribusi dari data dan menggambar hasilnya dengan \textgreater{}bar
pada plot3d, atau dengan plot kolom.
\end{eulercomment}
\begin{eulerprompt}
>w=normal(1000); // distribusi normal 0-1
>\{x,y\}=histo(w,10,v=[-6,-4,-2,-1,0,1,2,4,6]); // batas-batas interval v
>plot2d(x,y,>bar):
\end{eulerprompt}
\eulerimg{21}{images/EMT4Plot2D_Indonesia-114.png}
\begin{eulercomment}
Fungsi statplot() mengatur gaya dengan string sederhana.
\end{eulercomment}
\begin{eulerprompt}
>statplot(1:10,cumsum(random(10)),"b"):
\end{eulerprompt}
\eulerimg{21}{images/EMT4Plot2D_Indonesia-115.png}
\begin{eulerprompt}
>n=10; i=0:n; ...
>plot2d(i,bin(n,i)/2^n,a=0,b=10,c=0,d=0.3); ...
>plot2d(i,bin(n,i)/2^n,points=true,style="ow",add=true,color=blue):
\end{eulerprompt}
\eulerimg{21}{images/EMT4Plot2D_Indonesia-116.png}
\begin{eulercomment}
Selain itu, data dapat digambar sebagai batang. Dalam hal ini, x harus terurut dan
satu elemen lebih panjang daripada y. Batang akan membentang dari x[i] sampai x[i+1]
dengan nilai y[i]. Jika x berukuran sama dengan y, x akan diperpanjang satu elemen
dengan jarak terakhir.

Gaya isian dapat digunakan seperti di atas.
\end{eulercomment}
\begin{eulerprompt}
>n=10; k=bin(n,0:n); ...
>plot2d(-0.5:n+0.5,k,bar=true,fillcolor=lightgray):
\end{eulerprompt}
\eulerimg{21}{images/EMT4Plot2D_Indonesia-117.png}
\begin{eulercomment}
Data untuk plot batang (bar=1) dan histogram (histogram=1) dapat diberikan secara
eksplisit pada xv dan yv, atau dapat dihitung dari distribusi empiris pada xv dengan
\textgreater{}distribution (atau distribution=n). Histogram nilai-nilai xv akan dihitung otomatis
dengan \textgreater{}histogram. Jika \textgreater{}even ditentukan, nilai-nilai xv akan dihitung pada interval
bilangan bulat.
\end{eulercomment}
\begin{eulerprompt}
>plot2d(normal(10000),distribution=50):
\end{eulerprompt}
\eulerimg{21}{images/EMT4Plot2D_Indonesia-118.png}
\begin{eulerprompt}
>k=0:10; m=bin(10,k); x=(0:11)-0.5; plot2d(x,m,>bar):
\end{eulerprompt}
\eulerimg{21}{images/EMT4Plot2D_Indonesia-119.png}
\begin{eulerprompt}
>columnsplot(m,k):
\end{eulerprompt}
\eulerimg{21}{images/EMT4Plot2D_Indonesia-120.png}
\begin{eulerprompt}
>plot2d(random(600)*6,histogram=6):
\end{eulerprompt}
\eulerimg{21}{images/EMT4Plot2D_Indonesia-121.png}
\begin{eulercomment}
Untuk distribusi, ada parameter distribution=n, yang menghitung nilai secara otomatis
dan mencetak distribusi relatif dengan n subinterval.
\end{eulercomment}
\begin{eulerprompt}
>plot2d(normal(1,1000),distribution=10,style="\(\backslash\)/"):
\end{eulerprompt}
\eulerimg{21}{images/EMT4Plot2D_Indonesia-122.png}
\begin{eulercomment}
Dengan parameter even=true, interval bilangan bulat akan digunakan.
\end{eulercomment}
\begin{eulerprompt}
>plot2d(intrandom(1,1000,10),distribution=10,even=true):
\end{eulerprompt}
\eulerimg{21}{images/EMT4Plot2D_Indonesia-123.png}
\begin{eulercomment}
Perhatikan bahwa ada banyak plot statistik yang mungkin berguna. Lihat tutorial
tentang statistika.
\end{eulercomment}
\begin{eulerprompt}
>columnsplot(getmultiplicities(1:6,intrandom(1,6000,6))):
\end{eulerprompt}
\eulerimg{21}{images/EMT4Plot2D_Indonesia-124.png}
\begin{eulerprompt}
>plot2d(normal(1,1000),>distribution); ...
>  plot2d("qnormal(x)",color=red,thickness=2,>add):
\end{eulerprompt}
\eulerimg{21}{images/EMT4Plot2D_Indonesia-125.png}
\begin{eulercomment}
Ada pula banyak plot khusus untuk statistika. Boxplot menunjukkan kuartil dari
distribusi ini dan banyak pencilan (outlier). Menurut definisi, pencilan pada boxplot
adalah data yang melebihi 1,5 kali rentang 50\% tengah plot.
\end{eulercomment}
\begin{eulerprompt}
>M=normal(5,1000); boxplot(quartiles(M)):
\end{eulerprompt}
\eulerimg{21}{images/EMT4Plot2D_Indonesia-126.png}
\eulerheading{Fungsi Implisit}
\begin{eulercomment}
Plot implisit menunjukkan garis-garis level yang menyelesaikan f(x,y)=level, di mana
"level" dapat berupa satu nilai atau vektor nilai. Jika level="auto", akan ada nc
garis level, yang tersebar merata antara nilai minimum dan maksimum fungsi. Warna yang
lebih gelap atau lebih terang dapat ditambahkan dengan \textgreater{}hue untuk menunjukkan nilai
fungsi. Untuk fungsi implisit, xv harus berupa fungsi atau ekspresi dalam parameter x
dan y, atau alternatifnya, xv dapat berupa matriks nilai.

Euler dapat menandai garis-garis level

\end{eulercomment}
\begin{eulerformula}
\[
f(x,y) = c
\]
\end{eulerformula}
\begin{eulercomment}
dari fungsi apa pun.

Untuk menggambar himpunan f(x,y)=c untuk satu atau lebih konstanta c, Anda dapat
menggunakan plot2d() dengan plot implisitnya pada bidang. Parameter untuk c adalah
level=c, di mana c dapat berupa vektor garis level. Selain itu, skema warna dapat
digambar di latar belakang untuk menunjukkan nilai fungsi pada setiap titik di plot.
Parameter "n" menentukan kehalusan plot.
\end{eulercomment}
\begin{eulerprompt}
>aspect(1.5); 
>plot2d("x^2+y^2-x*y-x",r=1.5,level=0,contourcolor=red):
\end{eulerprompt}
\eulerimg{17}{images/EMT4Plot2D_Indonesia-127.png}
\begin{eulerprompt}
>expr := "2*x^2+x*y+3*y^4+y"; // definisikan ekspresi f(x,y)
>plot2d(expr,level=0): // Solusi f(x,y)=0
\end{eulerprompt}
\eulerimg{17}{images/EMT4Plot2D_Indonesia-128.png}
\begin{eulerprompt}
>plot2d(expr,level=0:0.5:20,>hue,contourcolor=white,n=200): // bagus
\end{eulerprompt}
\eulerimg{17}{images/EMT4Plot2D_Indonesia-129.png}
\begin{eulerprompt}
>plot2d(expr,level=0:0.5:20,>hue,>spectral,n=200,grid=4): // lebih bagus
\end{eulerprompt}
\eulerimg{17}{images/EMT4Plot2D_Indonesia-130.png}
\begin{eulercomment}
Ini juga berlaku untuk plot data. Tetapi Anda harus menentukan rentang\\
untuk label sumbu.
\end{eulercomment}
\begin{eulerprompt}
>x=-2:0.05:1; y=x'; z=expr(x,y);
>plot2d(z,level=0,a=-1,b=2,c=-2,d=1,>hue):
\end{eulerprompt}
\eulerimg{17}{images/EMT4Plot2D_Indonesia-131.png}
\begin{eulerprompt}
>plot2d("x^3-y^2",>contour,>hue,>spectral):
\end{eulerprompt}
\eulerimg{17}{images/EMT4Plot2D_Indonesia-132.png}
\begin{eulerprompt}
>plot2d("x^3-y^2",level=0,contourwidth=3,>add,contourcolor=red):
\end{eulerprompt}
\eulerimg{17}{images/EMT4Plot2D_Indonesia-133.png}
\begin{eulerprompt}
>z=z+normal(size(z))*0.2;
>plot2d(z,level=0.5,a=-1,b=2,c=-2,d=1):
\end{eulerprompt}
\eulerimg{17}{images/EMT4Plot2D_Indonesia-134.png}
\begin{eulerprompt}
>plot2d(expr,level=[0:0.2:5;0.05:0.2:5.05],color=lightgray):
\end{eulerprompt}
\eulerimg{17}{images/EMT4Plot2D_Indonesia-135.png}
\begin{eulerprompt}
>plot2d("x^2+y^3+x*y",level=1,r=4,n=100):
\end{eulerprompt}
\eulerimg{17}{images/EMT4Plot2D_Indonesia-136.png}
\begin{eulerprompt}
>plot2d("x^2+2*y^2-x*y",level=0:0.1:10,n=100,contourcolor=white,>hue):
\end{eulerprompt}
\eulerimg{17}{images/EMT4Plot2D_Indonesia-137.png}
\begin{eulercomment}
Dimungkinkan juga untuk mengisi himpunan

\end{eulercomment}
\begin{eulerformula}
\[
a \le f(x,y) \le b
\]
\end{eulerformula}
\begin{eulercomment}
dengan rentang level.

Dimungkinkan mengisi daerah nilai untuk fungsi tertentu. Untuk ini, level harus berupa
matriks 2xn. Baris pertama adalah batas bawah dan baris kedua berisi batas atas.
\end{eulercomment}
\begin{eulerprompt}
>plot2d(expr,level=[0;1],style="-",color=blue): // 0 <= f(x,y) <= 1
\end{eulerprompt}
\eulerimg{17}{images/EMT4Plot2D_Indonesia-138.png}
\begin{eulercomment}
Plot implisit juga dapat menampilkan rentang level. Dalam hal ini level harus berupa
matriks 2xn berisi interval level, di mana baris pertama berisi awal dan baris kedua
berisi akhir setiap interval. Alternatifnya, vektor baris sederhana dapat digunakan
untuk level, dan parameter dl memperluas nilai level menjadi interval.
\end{eulercomment}
\begin{eulerprompt}
>plot2d("x^4+y^4",r=1.5,level=[0;1],color=blue,style="/"):
\end{eulerprompt}
\eulerimg{17}{images/EMT4Plot2D_Indonesia-139.png}
\begin{eulerprompt}
>plot2d("x^2+y^3+x*y",level=[0,2,4;1,3,5],style="/",r=2,n=100):
\end{eulerprompt}
\eulerimg{17}{images/EMT4Plot2D_Indonesia-140.png}
\begin{eulerprompt}
>plot2d("x^2+y^3+x*y",level=-10:20,r=2,style="-",dl=0.1,n=100):
\end{eulerprompt}
\eulerimg{17}{images/EMT4Plot2D_Indonesia-141.png}
\begin{eulerprompt}
>plot2d("sin(x)*cos(y)",r=pi,>hue,>levels,n=100):
\end{eulerprompt}
\eulerimg{17}{images/EMT4Plot2D_Indonesia-142.png}
\begin{eulercomment}
Dimungkinkan juga menandai suatu daerah

\end{eulercomment}
\begin{eulerformula}
\[
a \le f(x,y) \le b.
\]
\end{eulerformula}
\begin{eulercomment}
Ini dilakukan dengan menambahkan level dengan dua baris.
\end{eulercomment}
\begin{eulerprompt}
>plot2d("(x^2+y^2-1)^3-x^2*y^3",r=1.3, ...
>  style="#",color=red,<outline, ...
>  level=[-2;0],n=100):
\end{eulerprompt}
\eulerimg{17}{images/EMT4Plot2D_Indonesia-143.png}
\begin{eulercomment}
Dimungkinkan menentukan level tertentu. Misalnya, kita dapat menggambar solusi
persamaan seperti

\end{eulercomment}
\begin{eulerformula}
\[
x^3-xy+x^2y^2=6
\]
\end{eulerformula}
\begin{eulerprompt}
>plot2d("x^3-x*y+x^2*y^2",r=6,level=1,n=100):
\end{eulerprompt}
\eulerimg{17}{images/EMT4Plot2D_Indonesia-144.png}
\begin{eulerprompt}
>function starplot1 (v, style="/", color=green, lab=none) ...
\end{eulerprompt}
\begin{eulerudf}
    if !holding() then clg; endif;
    w=window(); window(0,0,1024,1024);
    h=holding(1);
    r=max(abs(v))*1.2;
    setplot(-r,r,-r,r);
    n=cols(v); t=linspace(0,2pi,n);
    v=v|v[1]; c=v*cos(t); s=v*sin(t);
    cl=barcolor(color); st=barstyle(style);
    loop 1 to n
      polygon([0,c[#],c[#+1]],[0,s[#],s[#+1]],1);
      if lab!=none then
        rlab=v[#]+r*0.1;
        \{col,row\}=toscreen(cos(t[#])*rlab,sin(t[#])*rlab);
        ctext(""+lab[#],col,row-textheight()/2);
      endif;
    end;
    barcolor(cl); barstyle(st);
    holding(h);
    window(w);
  endfunction
\end{eulerudf}
\begin{eulercomment}
Di sini tidak ada grid atau tanda sumbu. Selain itu, kita menggunakan seluruh jendela
untuk plot.

Kita memanggil reset sebelum menguji plot ini untuk memulihkan pengaturan default
grafik. Hal ini tidak diperlukan jika Anda yakin plot Anda berfungsi.
\end{eulercomment}
\begin{eulerprompt}
>reset; starplot1(normal(1,10)+5,color=red,lab=1:10):
\end{eulerprompt}
\eulerimg{13}{images/EMT4Plot2D_Indonesia-145.png}
\begin{eulercomment}
Terkadang, Anda mungkin ingin menggambar sesuatu yang tidak dapat dilakukan plot2d,
tetapi hampir bisa.

Pada fungsi berikut, kita membuat plot impuls logaritmik. plot2d dapat membuat plot
logaritmik, tetapi tidak untuk batang impuls.
\end{eulercomment}
\begin{eulerprompt}
>function logimpulseplot1 (x,y) ...
\end{eulerprompt}
\begin{eulerudf}
    \{x0,y0\}=makeimpulse(x,log(y)/log(10));
    plot2d(x0,y0,>bar,grid=0);
    h=holding(1);
    frame();
    xgrid(ticks(x));
    p=plot();
    for i=-10 to 10;
      if i<=p[4] and i>=p[3] then
         ygrid(i,yt="10^"+i);
      endif;
    end;
    holding(h);
  endfunction
\end{eulerudf}
\begin{eulercomment}
Mari kita uji dengan nilai-nilai berdistribusi eksponensial.
\end{eulercomment}
\begin{eulerprompt}
>aspect(1.5); x=1:10; y=-log(random(size(x)))*200; ...
>logimpulseplot1(x,y):
\end{eulerprompt}
\eulerimg{17}{images/EMT4Plot2D_Indonesia-146.png}
\begin{eulercomment}
Mari kita animasikan kurva 2D menggunakan plot langsung. Perintah plot(x,y) cukup
menggambar kurva ke jendela plot. setplot(a,b,c,d) mengatur jendela ini.

Fungsi wait(0) memaksa plot muncul pada jendela grafik. Jika tidak, penggambaran ulang
hanya terjadi pada selang waktu yang jarang.
\end{eulercomment}
\begin{eulerprompt}
>function animliss (n,m) ...
\end{eulerprompt}
\begin{eulerudf}
  t=linspace(0,2pi,500);
  f=0;
  c=framecolor(0);
  l=linewidth(2);
  setplot(-1,1,-1,1);
  repeat
    clg;
    plot(sin(n*t),cos(m*t+f));
    wait(0);
    if testkey() then break; endif;
    f=f+0.02;
  end;
  framecolor(c);
  linewidth(l);
  endfunction
\end{eulerudf}
\begin{eulercomment}
Tekan sembarang tombol untuk menghentikan animasi ini.
\end{eulercomment}
\begin{eulerprompt}
>animliss(2,3); // lihat hasilnya, jika sudah puas, tekan ENTER
\end{eulerprompt}
\eulerheading{Plot Logaritmik}
\begin{eulercomment}
EMT menggunakan parameter "logplot" untuk skala logaritmik.\\
Plot logaritmik dapat digambar dengan skala logaritmik pada y menggunakan logplot=1,
atau skala logaritmik pada x dan y dengan logplot=2, atau pada x dengan logplot=3.

\end{eulercomment}
\begin{eulerttcomment}
 - logplot=1: y-logaritmik
 - logplot=2: x-y-logaritmik
 - logplot=3: x-logaritmik
\end{eulerttcomment}
\begin{eulerprompt}
>plot2d("exp(x^3-x)*x^2",1,5,logplot=1):
\end{eulerprompt}
\eulerimg{17}{images/EMT4Plot2D_Indonesia-147.png}
\begin{eulerprompt}
>plot2d("exp(x+sin(x))",0,100,logplot=1):
\end{eulerprompt}
\eulerimg{17}{images/EMT4Plot2D_Indonesia-148.png}
\begin{eulerprompt}
>plot2d("exp(x+sin(x))",10,100,logplot=2):
\end{eulerprompt}
\eulerimg{17}{images/EMT4Plot2D_Indonesia-149.png}
\begin{eulerprompt}
>plot2d("gamma(x)",1,10,logplot=1):
\end{eulerprompt}
\eulerimg{17}{images/EMT4Plot2D_Indonesia-150.png}
\begin{eulerprompt}
>plot2d("log(x*(2+sin(x/100)))",10,1000,logplot=3):
\end{eulerprompt}
\eulerimg{17}{images/EMT4Plot2D_Indonesia-151.png}
\begin{eulercomment}
Ini juga berlaku untuk plot data.
\end{eulercomment}
\begin{eulerprompt}
>x=10^(1:20); y=x^2-x;
>plot2d(x,y,logplot=2):
\end{eulerprompt}
\eulerimg{17}{images/EMT4Plot2D_Indonesia-152.png}
\eulersubheading{Contoh penyelesaian masalah menggambar kurva/plot 2D}
\begin{eulercomment}
Soal 1: Trigonometri \& Gelombang Terredam (Damped Oscillation)

Soal:Dalam bidang fisika dan teknik, fenomena ayunan yang meredam
seiring berjalannya waktu dimodelkan oleh fungsi peredaman. Gambarkan
grafik fungsi\\
\end{eulercomment}
\begin{eulerformula}
\[
f(x) = e^{-0.5x} \sin(3x) ; [0, 2\pi]
\]
\end{eulerformula}
\begin{eulercomment}
dengan\\
ketentuan:Buat kurva utama berwarna biru dengan ketebalan

2.Gambarkan pula garis kurva selubung (envelope)\\
\end{eulercomment}
\begin{eulerformula}
\[
y = e^{-0.5x} ; y = -e^{-0.5x}
\]
\end{eulerformula}
\begin{eulercomment}
menggunakan garis putus-putus berwarna merah sebagai batas peredaman.

3.Sertakan judul grafik, label sumbu x dan y, serta labelbox
keterangan kurva.
\end{eulercomment}
\begin{eulerprompt}
>reset;
>plot2d("exp(-0.5*x)*sin(3*x)", 0, 2*pi, color=blue, thickness=2, xl="x", yl="f(x)", title="Damped Oscillation"):
\end{eulerprompt}
\eulerimg{13}{images/EMT4Plot2D_Indonesia-155.png}
\begin{eulerprompt}
>plot2d("exp(-0.5*x)", >add, color=red, style="--"):
\end{eulerprompt}
\eulerimg{13}{images/EMT4Plot2D_Indonesia-156.png}
\begin{eulerprompt}
>plot2d("-exp(-0.5*x)", >add, color=red, style="--"):
\end{eulerprompt}
\eulerimg{13}{images/EMT4Plot2D_Indonesia-157.png}
\begin{eulerprompt}
>labelbox(["f(x) = e^\{-0.5x\} sin(3x)", "Envelope e^\{-0.5x\}"], colors=[blue, red], styles=["-", "--"]):
\end{eulerprompt}
\eulerimg{13}{images/EMT4Plot2D_Indonesia-158.png}
\begin{eulercomment}
Penjelasan Kode:plot2d\\
\end{eulercomment}
\begin{eulerformula}
\[
("exp(-0.5*x)*sin(3*x)", 0, 2*pi, ...)
\]
\end{eulerformula}
\begin{eulercomment}
menggambar fungsi gelombang utama pada domain 0 hingga 2\textbackslash{}pi.Parameter
\textgreater{}add digunakan untuk menimpa/menambahkan grafik selubung\\
\end{eulercomment}
\begin{eulerformula}
\[
(e^{-0.5x} dan -e^{-0.5x})
\]
\end{eulerformula}
\begin{eulercomment}
pada bidang koordinat yang sama tanpa menghapus grafik
sebelumnya.style="--" mengubah tampilan garis dari utuh menjadi garis
putus-putus.labelbox() membuat kotak legenda di pojok kanan atas untuk
membedakan fungsi utama dan garis batasnya.


Soal 2: Kalkulus – Kurva Parametrik (Lissajous Curve)\textbackslash{}\\
Soal:Suatu partikel bergerak di bidang 2D dengan posisi \\
\end{eulercomment}
\begin{eulerformula}
\[
(x(t), y(t))
\]
\end{eulerformula}
\begin{eulercomment}
yang terdefinisi secara parametrik oleh:\\
\end{eulercomment}
\begin{eulerformula}
\[
x(t) = \cos(3t), \quad y(t) = \sin(2t), \quad t \in [0, 2\pi]
\]
\end{eulerformula}
\begin{eulercomment}
Gambarkan lintasan kurva parametrik tersebut dalam bidang persegi
(square aspect ratio) dengan warna hijau dan isilah daerah dalam kurva
tersebut menggunakan arsir (filled/hatched style).
\end{eulercomment}
\begin{eulerprompt}
>reset;
>t = linspace(0, 2*pi, 1000);
>x = cos(3*t); y = sin(2*t);
>plot2d(x, y, >square, >filled, style="/", fillcolor=green, title="Lissajous Curve (3:2)"):
\end{eulerprompt}
\eulerimg{13}{images/EMT4Plot2D_Indonesia-163.png}
\begin{eulercomment}
Penjelasan Kode:linspace(0, 2*pi, 1000) membangkitkan 1000 titik
sampel halus untuk variabel parameter t.x = cos(3*t); y = sin(2*t);
menghitung posisi x dan y secara kualitatif berdasarkan parameter
t.Parameter \textgreater{}square memastikan rasio sumbu \textdollar{}X\textdollar{} dan \textdollar{}Y\textdollar{} seimbang (1:1)
agar bentuk geometris tidak terdistorsi.Parameter \textgreater{}filled beserta
style="/" mengisi bagian dalam loop kurva tertutup dengan pola
garis-garis diagonal berwarna hijau (fillcolor=green).

Soal 3: Aljabar – Kurva Implisit (Lemniscate of
Bernoulli)Soal:Gambarkan grafik persamaan implisit \\
\end{eulercomment}
\begin{eulerformula}
\[
f(x,y) = (x^2 + y^2)^2 - 4(x^2 - y^2) = 0
\]
\end{eulerformula}
\begin{eulercomment}
pada rentang \\
\end{eulercomment}
\begin{eulerformula}
\[
x \in [-2.5, 2.5]; y \in [-1.5, 1.5]. 
\]
\end{eulerformula}
\begin{eulercomment}
Buatlah agar kontur bernilai nol (f(x,y) = 0) terlihat tebal berwarna
merah dan tambahkan pewarnaan gradien (hue) di sekitar kontur untuk
menunjukkan perubahan nilai fungsi.
\end{eulercomment}
\begin{eulerprompt}
>reset;
>expr := "(x^2 + y^2)^2 - 4*(x^2 - y^2)";
>plot2d(expr, a=-2.5, b=2.5, c=-1.5, d=1.5, level=0:0.5:5, >hue, >spectral, n=200):
\end{eulerprompt}
\eulerimg{13}{images/EMT4Plot2D_Indonesia-166.png}
\begin{eulerprompt}
>plot2d(expr, level=0, contourcolor=red, contourwidth=3, >add):
\end{eulerprompt}
\eulerimg{13}{images/EMT4Plot2D_Indonesia-167.png}
\begin{eulercomment}
Penjelasan Kode:expr := ... mendefinisikan persamaan implisit sebagai
string ekspresi matematika dua variabel.a=-2.5, b=2.5, c=-1.5, d=1.5
menetapkan batas spesifik sumbu X (a ke b) dan sumbu Y (c ke
d).level=0:0.5:5 dipadukan dengan \textgreater{}hue dan \textgreater{}spectral memberikan warna
latar belakang (heatmap/gradient) berdasarkan spektrum nilai fungsi
f(x,y).Pemanggilan plot2d kedua dengan level=0, contourcolor=red, dan
contourwidth=3 mempertegas garis solusi persis pada f(x,y) = 0 dengan
warna merah tebal di atas latar spektrum.

Soal 4: Kalkulus – Deret Taylor \& Pengepasan FungsiSoal:Hampiri fungsi
\end{eulercomment}
\begin{eulerformula}
\[
f(x) = \cos(x)
\]
\end{eulerformula}
\begin{eulercomment}
menggunakan Polinomial Taylor orde ke-2, ke-4, dan ke-6 di sekitar x =
0. Gambarkan fungsi asli beserta ketiga polinomial hampirannya pada
interval \\
\end{eulercomment}
\begin{eulerformula}
\[
x \in [-\pi, \pi]
\]
\end{eulerformula}
\begin{eulercomment}
dengan ketentuan:Gunakan warna dan gaya garis yang berbeda untuk
setiap orde deret.Sertakan rumus LaTeX pada label grafik dan buat
sumbu koordinat dengan tanda tick khusus \\
\end{eulercomment}
\begin{eulerformula}
\[
(-\pi, 0, \pi).
\]
\end{eulerformula}
\begin{eulerprompt}
>reset;
>plot2d("cos(x)", -pi, pi, color=black, thickness=2, grid=4, <ticks):
\end{eulerprompt}
\eulerimg{13}{images/EMT4Plot2D_Indonesia-171.png}
\begin{eulerprompt}
>plot2d("1 - x^2/2", >add, color=blue, style="--"):
\end{eulerprompt}
\eulerimg{13}{images/EMT4Plot2D_Indonesia-172.png}
\begin{eulerprompt}
>plot2d("1 - x^2/2 + x^4/24", >add, color=red, style="-."):
\end{eulerprompt}
\eulerimg{13}{images/EMT4Plot2D_Indonesia-173.png}
\begin{eulerprompt}
>plot2d("1 - x^2/2 + x^4/24 - x^6/720", >add, color=green, style="."):
\end{eulerprompt}
\eulerimg{13}{images/EMT4Plot2D_Indonesia-174.png}
\begin{eulerprompt}
>ygrid(-1.5:0.5:1.5);
>xgrid([-pi, 0, pi]);
>xtick([-pi, 0, pi], ["-\(\backslash\)pi", "0", "\(\backslash\)pi"], >latex);
>labelbox(["cos(x)", "T2", "T4", "T6"], colors=[black, blue, red, green], styles=["-", "--", "-.", "."]):
\end{eulerprompt}
\eulerimg{13}{images/EMT4Plot2D_Indonesia-175.png}
\begin{eulercomment}
Penjelasan Kode:grid=4 dan \textless{}ticks mematikan garis kisi dan tanda sumbu
bawaan agar kita dapat membuat sumbu kustom.xgrid(...) dan ygrid(...)
menambahkan garis kisi secara manual pada lokasi yang
ditentukan.xtick([-pi, 0, pi], ["-\textbackslash{}pi", "0", "\textbackslash{}pi"], \textgreater{}latex) mengubah
label numerik biasa pada sumbu  menjadi simbol LaTeX =.Setiap
polinomial Taylor diplot bertumpuk dengan gaya garis berbeda (--, -*,
.) untuk memvisualisasikan bagaimana orde yang lebih tinggi memberikan
akurasi hampiran yang lebih baik.
\end{eulercomment}
\eulerheading{Rujukan Lengkap Fungsi plot2d()}
\begin{eulercomment}
\end{eulercomment}
\begin{eulerttcomment}
  function plot2d (xv, yv, btest, a, b, c, d, xmin, xmax, r, n,  ..
  logplot, grid, frame, framecolor, square, color, thickness, style, ..
  auto, add, user, delta, points, addpoints, pointstyle, bar, histogram,  ..
  distribution, even, steps, own, adaptive, hue, level, contour,  ..
  nc, filled, fillcolor, outline, title, xl, yl, maps, contourcolor, ..
  contourwidth, ticks, margin, clipping, cx, cy, insimg, spectral,  ..
  cgrid, vertical, smaller, dl, niveau, levels)
\end{eulerttcomment}
\begin{eulercomment}
Fungsi plot serbaguna untuk plot pada bidang (plot 2D). Fungsi ini dapat membuat
plot fungsi satu variabel, plot data, kurva pada bidang, plot batang, grid bilangan
kompleks, dan plot implisit fungsi dua variabel.

Parameter
\\
x,y       : persamaan, fungsi, atau vektor data\\
a,b,c,d   : Area plot (default a=-2,b=2)\\
r         : jika r diatur, maka a=cx-r, b=cx+r, c=cy-r, d=cy+r\\
\end{eulercomment}
\begin{eulerttcomment}
            r dapat berupa vektor [rx,ry] atau vektor [rx1,rx2,ry1,ry2].
\end{eulerttcomment}
\begin{eulercomment}
xmin,xmax : rentang parameter untuk kurva\\
auto      : Tentukan rentang y secara otomatis (default)\\
square    : jika true, usahakan menjaga rentang x-y tetap persegi\\
n         : banyak interval (default adaptif)\\
grid      : 0 = tanpa grid dan label,\\
\end{eulercomment}
\begin{eulerttcomment}
            1 = hanya sumbu,
            2 = grid normal (lihat di bawah untuk banyaknya garis grid)
            3 = sumbu di dalam
            4 = tanpa grid
            5 = grid penuh termasuk margin
            6 = tick pada bingkai
            7 = hanya sumbu
            8 = hanya sumbu, sub-tick
\end{eulerttcomment}
\begin{eulercomment}
frame     : 0 = tanpa bingkai\\
framecolor: warna bingkai dan grid\\
margin    : bilangan antara 0 dan 0.4 untuk margin di sekitar plot\\
color     : Warna kurva. Jika berupa vektor warna,\\
\end{eulercomment}
\begin{eulerttcomment}
            warna tersebut akan digunakan untuk setiap baris matriks plot. Pada
            plot titik, harus berupa vektor kolom. Jika vektor baris atau
            matriks warna penuh digunakan untuk plot titik, warna tersebut akan
            digunakan untuk setiap titik data.
\end{eulerttcomment}
\begin{eulercomment}
thickness : ketebalan garis untuk kurva\\
\end{eulercomment}
\begin{eulerttcomment}
            Nilai ini boleh lebih kecil dari 1 untuk garis yang sangat tipis.
\end{eulerttcomment}
\begin{eulercomment}
style     : Gaya plot untuk garis, penanda, dan isian.\\
\end{eulercomment}
\begin{eulerttcomment}
            Untuk titik gunakan
            "[]", "<>", ".", "..", "...",
            "*", "+", "|", "-", "o"
            "[]#", "<>#", "o#" (bentuk terisi)
            "[]w", "<>w", "ow" (tidak transparan)
            Untuk garis gunakan
            "-", "--", "-.", ".", ".-.", "-.-", "->"
            Untuk poligon terisi atau plot batang gunakan
            "#", "#O", "O", "/", "\(\backslash\)", "\(\backslash\)/",
            "+", "|", "-", "t"
\end{eulerttcomment}
\begin{eulercomment}
points    : gambar titik tunggal sebagai ganti ruas garis\\
addpoints : jika true, gambar ruas garis dan titik\\
add       : tambahkan plot ke plot yang sudah ada\\
user      : aktifkan interaksi pengguna untuk fungsi\\
delta     : ukuran langkah untuk interaksi pengguna\\
bar       : plot batang (x adalah batas interval, y nilai interval)\\
histogram : menggambar frekuensi x pada n subinterval\\
distribution=n : menggambar distribusi x dengan n subinterval\\
even      : gunakan nilai bulat untuk histogram otomatis.\\
steps     : menggambar fungsi sebagai fungsi tangga (steps=1,2)\\
adaptive  : gunakan plot adaptif (n adalah banyak langkah minimal)\\
level     : menggambar garis level fungsi implisit dua variabel\\
outline   : menggambar batas rentang level.
\\
Jika nilai level berupa matriks 2xn, rentang level akan digambar\\
dengan warna menggunakan gaya isian yang diberikan. Jika outline true,\\
batasnya akan digambar dengan warna kontur. Dengan fitur ini, daerah\\
f(x,y) di antara batas-batas tertentu dapat ditandai.
\\
hue       : tambahkan warna hue pada plot level untuk menunjukkan nilai\\
\end{eulercomment}
\begin{eulerttcomment}
            fungsi
\end{eulerttcomment}
\begin{eulercomment}
contour   : Gunakan plot level dengan level otomatis\\
nc        : banyaknya garis level otomatis\\
title     : judul plot (default "")\\
xl, yl    : label untuk sumbu x dan y\\
smaller   : jika \textgreater{}0, akan ada lebih banyak ruang di kiri untuk label.\\
vertical  :\\
\end{eulercomment}
\begin{eulerttcomment}
  Menyalakan atau mematikan label vertikal. Ini mengubah variabel global
  verticallabels secara lokal untuk satu plot. Nilai 1 hanya mengatur teks
  vertikal, nilai 2 menggunakan label numerik vertikal pada sumbu y.
\end{eulerttcomment}
\begin{eulercomment}
filled    : isi plot sebuah kurva\\
fillcolor : warna isian untuk batang dan kurva terisi\\
outline   : batas untuk poligon terisi\\
logplot   : atur plot logaritmik\\
\end{eulercomment}
\begin{eulerttcomment}
            1 = logplot pada y,
            2 = logplot pada xy,
            3 = logplot pada x
\end{eulerttcomment}
\begin{eulercomment}
own       :\\
\end{eulercomment}
\begin{eulerttcomment}
  String yang menunjuk ke rutin plot buatan sendiri. Dengan >user, Anda
  mendapatkan interaksi pengguna yang sama seperti pada plot2d. Rentang akan
  diatur sebelum setiap pemanggilan fungsi Anda.
\end{eulerttcomment}
\begin{eulercomment}
maps      : petakan ekspresi (0 lebih cepat), fungsi selalu dipetakan.\\
contourcolor : warna garis kontur\\
contourwidth : lebar garis kontur\\
clipping  : mengaktifkan/menonaktifkan pemotongan (default true)\\
title     :\\
\end{eulercomment}
\begin{eulerttcomment}
  Ini dapat digunakan untuk mendeskripsikan plot. Judul akan muncul di atas
  plot. Selain itu, label untuk sumbu x dan y dapat ditambahkan dengan
  xl="string" atau yl="string". Label lain dapat ditambahkan dengan fungsi
  label() atau labelbox(). Judul dapat berupa string unicode atau gambar
  rumus Latex.
\end{eulerttcomment}
\begin{eulercomment}
cgrid     :\\
\end{eulercomment}
\begin{eulerttcomment}
  Menentukan banyaknya garis grid untuk plot grid kompleks.
  Sebaiknya merupakan pembagi dari ukuran matriks dikurangi 1 (banyaknya
  subinterval). cgrid dapat berupa vektor [cx,cy].
\end{eulerttcomment}
\begin{eulercomment}

Gambaran Umum

Fungsi ini dapat menggambar

- ekspresi, call collection, atau fungsi satu variabel,\\
- kurva parametrik,\\
- data x terhadap data y,\\
- fungsi implisit,\\
- plot batang,\\
- grid kompleks,\\
- poligon.

Jika fungsi atau ekspresi untuk xv diberikan, plot2d() akan menghitung\\
nilai pada rentang yang diberikan menggunakan fungsi atau ekspresi tersebut.\\
Ekspresi harus berupa ekspresi dalam variabel x. Rentang harus\\
didefinisikan pada parameter a dan b, kecuali rentang default\\
[-2,2] yang ingin digunakan. Rentang y akan dihitung otomatis,\\
kecuali c dan d ditentukan, atau jari-jari r, yang menghasilkan rentang\\
[-r,r] untuk x dan y. Untuk plot fungsi, plot2d secara default akan\\
menggunakan evaluasi fungsi adaptif. Untuk mempercepat\\
plot fungsi yang rumit, matikan dengan \textless{}adaptive, dan\\
jika perlu kurangi banyaknya interval n. Selain itu, plot2d()\\
secara default akan menggunakan pemetaan (mapping), yaitu menghitung elemen\\
plot satu per satu. Jika ekspresi atau fungsi Anda dapat menangani\\
vektor x, Anda dapat mematikannya dengan \textless{}maps untuk evaluasi yang lebih cepat.

Perhatikan bahwa plot adaptif selalu dihitung elemen demi elemen. \\
Jika fungsi atau ekspresi untuk xv dan untuk yv sama-sama ditentukan,\\
plot2d() akan menghitung kurva dengan nilai xv sebagai koordinat x\\
dan nilai yv sebagai koordinat y. Dalam hal ini, rentang harus\\
didefinisikan untuk parameter menggunakan xmin, xmax. Ekspresi yang termuat\\
dalam string harus selalu berupa ekspresi dalam variabel parameter x.
\end{eulercomment}
\end{eulernotebook}
\end{document}
